\documentclass[11pt,a4paper]{article}

\newcommand{\CourseCode}{CSAI2017P: Applied Machine Learning}
\newcommand{\AssignmentName}{Mid-Semester Question Bank -- Units I--V(pt)}

% Shared settings for Applied Machine Learning lab assignments and marking schemes.
% Layered on the house notes preamble; keep free of \documentclass.
%
% Callers must define \CourseCode, \AssignmentName before \begin{document}.

% Depth-agnostic locate of the house notes preamble: briefs compile from
% public/lab-assignments/LAxx/latex/ and marking schemes from
% private/solutions/lab-assignments/LAxx/latex/, which sit at different depths.
\IfFileExists{../../../../public/_shared/notes-common.tex}{%
  % Shared notes settings for Applied Machine Learning lectures (UPES).
% Keep this file free of \documentclass to allow reuse via \input.

\usepackage[utf8]{inputenc}
\usepackage[T1]{fontenc}
\usepackage{lmodern}          % scalable T1 fonts; without this pdflatex falls back
\input{glyphtounicode}        % to bitmapped EC Type 3 and the PDFs are neither
\pdfgentounicode=1            % searchable nor screen-reader friendly
\usepackage[a4paper,margin=1in]{geometry}
\usepackage{hyperref}
\usepackage{xurl}
\usepackage{booktabs}
\usepackage{amsmath}
\usepackage{amssymb}
\usepackage{listings}
\usepackage{xcolor}
\usepackage{enumitem}
\usepackage{parskip}

% Code listings (Python-oriented for this subject).
\lstset{
  language=Python,
  basicstyle=\ttfamily\small,
  keywordstyle=\color{blue},
  commentstyle=\color{gray},
  stringstyle=\color{teal},
  showstringspaces=false,
  columns=fullflexible,
  frame=single,
  framerule=0.2pt,
  breaklines=true
}

\setlist[itemize]{topsep=0.3em,itemsep=0.2em}
\setlist[enumerate]{topsep=0.3em,itemsep=0.2em}

% Simple per-section source line.
\newcommand{\notesources}[1]{\par\noindent{\small\color{gray}\textit{Sources:} \sloppy #1\par}}

% Unit-wise source strings for this subject (used by slides + notes).
% Path is relative to the lecture `latex/` folder (the compilation CWD).
\IfFileExists{../../../_shared/aml-sources.tex}{%
  % Source strings for Applied Machine Learning (CSAI2017P) — used by slides + notes.
% Prescribed texts and reference books are taken from the course syllabus.
% Support texts are the author-hosted open-access editions:
%   statlearning.com            (ISL with Python)
%   probml.github.io/pml-book   (Murphy, Probabilistic Machine Learning)
%   d2l.ai                      (Dive into Deep Learning)
%   mml-book.github.io          (Mathematics for Machine Learning)

\newcommand{\AMLRepoURL}{https://github.com/tali7c/Applied-Machine-Learning}

% --- prescribed -------------------------------------------------------------
\newcommand{\AMLPrescribed}{%
M\"uller and Guido, \textit{Introduction to Machine Learning with Python}, O'Reilly, 2016.;%
Bishop, \textit{Pattern Recognition and Machine Learning}, Springer, 2006.;%
Murphy, \textit{Machine Learning: A Probabilistic Perspective}, MIT Press, 2012.;%
Han, Kamber and Pei, \textit{Data Mining: Concepts and Techniques} (3e), Morgan Kaufmann, 2012.%
}

% --- unit-wise ---------------------------------------------------------------
\newcommand{\AMLUnitISources}{%
M\"uller and Guido, \textit{Introduction to Machine Learning with Python}, Ch. 1.;%
Bishop, \textit{Pattern Recognition and Machine Learning}, Ch. 1.;%
Murphy, \textit{Probabilistic Machine Learning: An Introduction}, MIT Press, 2022, Ch. 1.;%
Mitchell, \textit{Machine Learning}, McGraw Hill, 1997, Ch. 1.;%
Scikit-learn documentation, accessed 2026-08-04.%
}

\newcommand{\AMLUnitIISources}{%
Bishop, \textit{Pattern Recognition and Machine Learning}, \S1.5, \S4.3, \S7.1.;%
Zhang, Lipton, Li and Smola, \textit{Dive into Deep Learning}, \S3.4.;%
Shalev-Shwartz and Ben-David, \textit{Understanding Machine Learning}, Ch. 15.;%
Murphy, \textit{Probabilistic Machine Learning: An Introduction}, Ch. 5.%
}

\newcommand{\AMLUnitIIISources}{%
Zhang, Lipton, Li and Smola, \textit{Dive into Deep Learning}, Ch. 12 (Optimization Algorithms).;%
Deisenroth, Faisal and Ong, \textit{Mathematics for Machine Learning}, Ch. 7.;%
Kingma and Ba, ``Adam: A Method for Stochastic Optimization,'' ICLR 2015.;%
Ruder, ``An overview of gradient descent optimization algorithms,'' arXiv:1609.04747, 2016.%
}

\newcommand{\AMLUnitIVSources}{%
Han, Kamber and Pei, \textit{Data Mining: Concepts and Techniques} (3e), Ch. 3.;%
M\"uller and Guido, \textit{Introduction to Machine Learning with Python}, Ch. 3--4.;%
James, Witten, Hastie, Tibshirani and Taylor, \textit{An Introduction to Statistical Learning with Python}, Ch. 5--6, 12.;%
Scikit-learn documentation (preprocessing, model selection), accessed 2026-08-04.%
}

\newcommand{\AMLUnitVSources}{%
James et al., \textit{An Introduction to Statistical Learning with Python}, Ch. 3--6.;%
Bishop, \textit{Pattern Recognition and Machine Learning}, Ch. 3.;%
Deisenroth, Faisal and Ong, \textit{Mathematics for Machine Learning}, Ch. 9.;%
M\"uller and Guido, \textit{Introduction to Machine Learning with Python}, \S2.3.%
}

\newcommand{\AMLUnitVISources}{%
Han, Kamber and Pei, \textit{Data Mining: Concepts and Techniques} (3e), Ch. 8--9.;%
James et al., \textit{An Introduction to Statistical Learning with Python}, Ch. 8--9.;%
Bishop, \textit{Pattern Recognition and Machine Learning}, Ch. 4--5, 7, 14.;%
Zhang et al., \textit{Dive into Deep Learning}, Ch. 5.%
}

\newcommand{\AMLUnitVIISources}{%
Han, Kamber and Pei, \textit{Data Mining: Concepts and Techniques} (3e), Ch. 10.;%
James et al., \textit{An Introduction to Statistical Learning with Python}, Ch. 12.;%
M\"uller and Guido, \textit{Introduction to Machine Learning with Python}, \S3.5.;%
Ester, Kriegel, Sander and Xu, ``A density-based algorithm for discovering clusters (DBSCAN),'' KDD 1996.%
}

\newcommand{\AMLLabSources}{%
M\"uller and Guido, \textit{Introduction to Machine Learning with Python}, companion notebooks.;%
McKinney, \textit{Python for Data Analysis} (3e), O'Reilly, 2022.;%
Scikit-learn user guide, accessed 2026-08-04.;%
Matplotlib and Seaborn documentation, accessed 2026-08-04.%
}
%
}{%
  % Faculty copies live under `private/` so their compile CWD is different.
  \IfFileExists{../../../../public/_shared/aml-sources.tex}{%
    % Source strings for Applied Machine Learning (CSAI2017P) — used by slides + notes.
% Prescribed texts and reference books are taken from the course syllabus.
% Support texts are the author-hosted open-access editions:
%   statlearning.com            (ISL with Python)
%   probml.github.io/pml-book   (Murphy, Probabilistic Machine Learning)
%   d2l.ai                      (Dive into Deep Learning)
%   mml-book.github.io          (Mathematics for Machine Learning)

\newcommand{\AMLRepoURL}{https://github.com/tali7c/Applied-Machine-Learning}

% --- prescribed -------------------------------------------------------------
\newcommand{\AMLPrescribed}{%
M\"uller and Guido, \textit{Introduction to Machine Learning with Python}, O'Reilly, 2016.;%
Bishop, \textit{Pattern Recognition and Machine Learning}, Springer, 2006.;%
Murphy, \textit{Machine Learning: A Probabilistic Perspective}, MIT Press, 2012.;%
Han, Kamber and Pei, \textit{Data Mining: Concepts and Techniques} (3e), Morgan Kaufmann, 2012.%
}

% --- unit-wise ---------------------------------------------------------------
\newcommand{\AMLUnitISources}{%
M\"uller and Guido, \textit{Introduction to Machine Learning with Python}, Ch. 1.;%
Bishop, \textit{Pattern Recognition and Machine Learning}, Ch. 1.;%
Murphy, \textit{Probabilistic Machine Learning: An Introduction}, MIT Press, 2022, Ch. 1.;%
Mitchell, \textit{Machine Learning}, McGraw Hill, 1997, Ch. 1.;%
Scikit-learn documentation, accessed 2026-08-04.%
}

\newcommand{\AMLUnitIISources}{%
Bishop, \textit{Pattern Recognition and Machine Learning}, \S1.5, \S4.3, \S7.1.;%
Zhang, Lipton, Li and Smola, \textit{Dive into Deep Learning}, \S3.4.;%
Shalev-Shwartz and Ben-David, \textit{Understanding Machine Learning}, Ch. 15.;%
Murphy, \textit{Probabilistic Machine Learning: An Introduction}, Ch. 5.%
}

\newcommand{\AMLUnitIIISources}{%
Zhang, Lipton, Li and Smola, \textit{Dive into Deep Learning}, Ch. 12 (Optimization Algorithms).;%
Deisenroth, Faisal and Ong, \textit{Mathematics for Machine Learning}, Ch. 7.;%
Kingma and Ba, ``Adam: A Method for Stochastic Optimization,'' ICLR 2015.;%
Ruder, ``An overview of gradient descent optimization algorithms,'' arXiv:1609.04747, 2016.%
}

\newcommand{\AMLUnitIVSources}{%
Han, Kamber and Pei, \textit{Data Mining: Concepts and Techniques} (3e), Ch. 3.;%
M\"uller and Guido, \textit{Introduction to Machine Learning with Python}, Ch. 3--4.;%
James, Witten, Hastie, Tibshirani and Taylor, \textit{An Introduction to Statistical Learning with Python}, Ch. 5--6, 12.;%
Scikit-learn documentation (preprocessing, model selection), accessed 2026-08-04.%
}

\newcommand{\AMLUnitVSources}{%
James et al., \textit{An Introduction to Statistical Learning with Python}, Ch. 3--6.;%
Bishop, \textit{Pattern Recognition and Machine Learning}, Ch. 3.;%
Deisenroth, Faisal and Ong, \textit{Mathematics for Machine Learning}, Ch. 9.;%
M\"uller and Guido, \textit{Introduction to Machine Learning with Python}, \S2.3.%
}

\newcommand{\AMLUnitVISources}{%
Han, Kamber and Pei, \textit{Data Mining: Concepts and Techniques} (3e), Ch. 8--9.;%
James et al., \textit{An Introduction to Statistical Learning with Python}, Ch. 8--9.;%
Bishop, \textit{Pattern Recognition and Machine Learning}, Ch. 4--5, 7, 14.;%
Zhang et al., \textit{Dive into Deep Learning}, Ch. 5.%
}

\newcommand{\AMLUnitVIISources}{%
Han, Kamber and Pei, \textit{Data Mining: Concepts and Techniques} (3e), Ch. 10.;%
James et al., \textit{An Introduction to Statistical Learning with Python}, Ch. 12.;%
M\"uller and Guido, \textit{Introduction to Machine Learning with Python}, \S3.5.;%
Ester, Kriegel, Sander and Xu, ``A density-based algorithm for discovering clusters (DBSCAN),'' KDD 1996.%
}

\newcommand{\AMLLabSources}{%
M\"uller and Guido, \textit{Introduction to Machine Learning with Python}, companion notebooks.;%
McKinney, \textit{Python for Data Analysis} (3e), O'Reilly, 2022.;%
Scikit-learn user guide, accessed 2026-08-04.;%
Matplotlib and Seaborn documentation, accessed 2026-08-04.%
}
%
  }{%
    % Question papers live at `private/<family>/<instrument>/`.
    \IfFileExists{../../../public/_shared/aml-sources.tex}{%
      % Source strings for Applied Machine Learning (CSAI2017P) — used by slides + notes.
% Prescribed texts and reference books are taken from the course syllabus.
% Support texts are the author-hosted open-access editions:
%   statlearning.com            (ISL with Python)
%   probml.github.io/pml-book   (Murphy, Probabilistic Machine Learning)
%   d2l.ai                      (Dive into Deep Learning)
%   mml-book.github.io          (Mathematics for Machine Learning)

\newcommand{\AMLRepoURL}{https://github.com/tali7c/Applied-Machine-Learning}

% --- prescribed -------------------------------------------------------------
\newcommand{\AMLPrescribed}{%
M\"uller and Guido, \textit{Introduction to Machine Learning with Python}, O'Reilly, 2016.;%
Bishop, \textit{Pattern Recognition and Machine Learning}, Springer, 2006.;%
Murphy, \textit{Machine Learning: A Probabilistic Perspective}, MIT Press, 2012.;%
Han, Kamber and Pei, \textit{Data Mining: Concepts and Techniques} (3e), Morgan Kaufmann, 2012.%
}

% --- unit-wise ---------------------------------------------------------------
\newcommand{\AMLUnitISources}{%
M\"uller and Guido, \textit{Introduction to Machine Learning with Python}, Ch. 1.;%
Bishop, \textit{Pattern Recognition and Machine Learning}, Ch. 1.;%
Murphy, \textit{Probabilistic Machine Learning: An Introduction}, MIT Press, 2022, Ch. 1.;%
Mitchell, \textit{Machine Learning}, McGraw Hill, 1997, Ch. 1.;%
Scikit-learn documentation, accessed 2026-08-04.%
}

\newcommand{\AMLUnitIISources}{%
Bishop, \textit{Pattern Recognition and Machine Learning}, \S1.5, \S4.3, \S7.1.;%
Zhang, Lipton, Li and Smola, \textit{Dive into Deep Learning}, \S3.4.;%
Shalev-Shwartz and Ben-David, \textit{Understanding Machine Learning}, Ch. 15.;%
Murphy, \textit{Probabilistic Machine Learning: An Introduction}, Ch. 5.%
}

\newcommand{\AMLUnitIIISources}{%
Zhang, Lipton, Li and Smola, \textit{Dive into Deep Learning}, Ch. 12 (Optimization Algorithms).;%
Deisenroth, Faisal and Ong, \textit{Mathematics for Machine Learning}, Ch. 7.;%
Kingma and Ba, ``Adam: A Method for Stochastic Optimization,'' ICLR 2015.;%
Ruder, ``An overview of gradient descent optimization algorithms,'' arXiv:1609.04747, 2016.%
}

\newcommand{\AMLUnitIVSources}{%
Han, Kamber and Pei, \textit{Data Mining: Concepts and Techniques} (3e), Ch. 3.;%
M\"uller and Guido, \textit{Introduction to Machine Learning with Python}, Ch. 3--4.;%
James, Witten, Hastie, Tibshirani and Taylor, \textit{An Introduction to Statistical Learning with Python}, Ch. 5--6, 12.;%
Scikit-learn documentation (preprocessing, model selection), accessed 2026-08-04.%
}

\newcommand{\AMLUnitVSources}{%
James et al., \textit{An Introduction to Statistical Learning with Python}, Ch. 3--6.;%
Bishop, \textit{Pattern Recognition and Machine Learning}, Ch. 3.;%
Deisenroth, Faisal and Ong, \textit{Mathematics for Machine Learning}, Ch. 9.;%
M\"uller and Guido, \textit{Introduction to Machine Learning with Python}, \S2.3.%
}

\newcommand{\AMLUnitVISources}{%
Han, Kamber and Pei, \textit{Data Mining: Concepts and Techniques} (3e), Ch. 8--9.;%
James et al., \textit{An Introduction to Statistical Learning with Python}, Ch. 8--9.;%
Bishop, \textit{Pattern Recognition and Machine Learning}, Ch. 4--5, 7, 14.;%
Zhang et al., \textit{Dive into Deep Learning}, Ch. 5.%
}

\newcommand{\AMLUnitVIISources}{%
Han, Kamber and Pei, \textit{Data Mining: Concepts and Techniques} (3e), Ch. 10.;%
James et al., \textit{An Introduction to Statistical Learning with Python}, Ch. 12.;%
M\"uller and Guido, \textit{Introduction to Machine Learning with Python}, \S3.5.;%
Ester, Kriegel, Sander and Xu, ``A density-based algorithm for discovering clusters (DBSCAN),'' KDD 1996.%
}

\newcommand{\AMLLabSources}{%
M\"uller and Guido, \textit{Introduction to Machine Learning with Python}, companion notebooks.;%
McKinney, \textit{Python for Data Analysis} (3e), O'Reilly, 2022.;%
Scikit-learn user guide, accessed 2026-08-04.;%
Matplotlib and Seaborn documentation, accessed 2026-08-04.%
}
%
    }{%
      % Marking schemes live at `private/solutions/<family>/<id>/latex/`.
      \IfFileExists{../../../../../public/_shared/aml-sources.tex}{%
        % Source strings for Applied Machine Learning (CSAI2017P) — used by slides + notes.
% Prescribed texts and reference books are taken from the course syllabus.
% Support texts are the author-hosted open-access editions:
%   statlearning.com            (ISL with Python)
%   probml.github.io/pml-book   (Murphy, Probabilistic Machine Learning)
%   d2l.ai                      (Dive into Deep Learning)
%   mml-book.github.io          (Mathematics for Machine Learning)

\newcommand{\AMLRepoURL}{https://github.com/tali7c/Applied-Machine-Learning}

% --- prescribed -------------------------------------------------------------
\newcommand{\AMLPrescribed}{%
M\"uller and Guido, \textit{Introduction to Machine Learning with Python}, O'Reilly, 2016.;%
Bishop, \textit{Pattern Recognition and Machine Learning}, Springer, 2006.;%
Murphy, \textit{Machine Learning: A Probabilistic Perspective}, MIT Press, 2012.;%
Han, Kamber and Pei, \textit{Data Mining: Concepts and Techniques} (3e), Morgan Kaufmann, 2012.%
}

% --- unit-wise ---------------------------------------------------------------
\newcommand{\AMLUnitISources}{%
M\"uller and Guido, \textit{Introduction to Machine Learning with Python}, Ch. 1.;%
Bishop, \textit{Pattern Recognition and Machine Learning}, Ch. 1.;%
Murphy, \textit{Probabilistic Machine Learning: An Introduction}, MIT Press, 2022, Ch. 1.;%
Mitchell, \textit{Machine Learning}, McGraw Hill, 1997, Ch. 1.;%
Scikit-learn documentation, accessed 2026-08-04.%
}

\newcommand{\AMLUnitIISources}{%
Bishop, \textit{Pattern Recognition and Machine Learning}, \S1.5, \S4.3, \S7.1.;%
Zhang, Lipton, Li and Smola, \textit{Dive into Deep Learning}, \S3.4.;%
Shalev-Shwartz and Ben-David, \textit{Understanding Machine Learning}, Ch. 15.;%
Murphy, \textit{Probabilistic Machine Learning: An Introduction}, Ch. 5.%
}

\newcommand{\AMLUnitIIISources}{%
Zhang, Lipton, Li and Smola, \textit{Dive into Deep Learning}, Ch. 12 (Optimization Algorithms).;%
Deisenroth, Faisal and Ong, \textit{Mathematics for Machine Learning}, Ch. 7.;%
Kingma and Ba, ``Adam: A Method for Stochastic Optimization,'' ICLR 2015.;%
Ruder, ``An overview of gradient descent optimization algorithms,'' arXiv:1609.04747, 2016.%
}

\newcommand{\AMLUnitIVSources}{%
Han, Kamber and Pei, \textit{Data Mining: Concepts and Techniques} (3e), Ch. 3.;%
M\"uller and Guido, \textit{Introduction to Machine Learning with Python}, Ch. 3--4.;%
James, Witten, Hastie, Tibshirani and Taylor, \textit{An Introduction to Statistical Learning with Python}, Ch. 5--6, 12.;%
Scikit-learn documentation (preprocessing, model selection), accessed 2026-08-04.%
}

\newcommand{\AMLUnitVSources}{%
James et al., \textit{An Introduction to Statistical Learning with Python}, Ch. 3--6.;%
Bishop, \textit{Pattern Recognition and Machine Learning}, Ch. 3.;%
Deisenroth, Faisal and Ong, \textit{Mathematics for Machine Learning}, Ch. 9.;%
M\"uller and Guido, \textit{Introduction to Machine Learning with Python}, \S2.3.%
}

\newcommand{\AMLUnitVISources}{%
Han, Kamber and Pei, \textit{Data Mining: Concepts and Techniques} (3e), Ch. 8--9.;%
James et al., \textit{An Introduction to Statistical Learning with Python}, Ch. 8--9.;%
Bishop, \textit{Pattern Recognition and Machine Learning}, Ch. 4--5, 7, 14.;%
Zhang et al., \textit{Dive into Deep Learning}, Ch. 5.%
}

\newcommand{\AMLUnitVIISources}{%
Han, Kamber and Pei, \textit{Data Mining: Concepts and Techniques} (3e), Ch. 10.;%
James et al., \textit{An Introduction to Statistical Learning with Python}, Ch. 12.;%
M\"uller and Guido, \textit{Introduction to Machine Learning with Python}, \S3.5.;%
Ester, Kriegel, Sander and Xu, ``A density-based algorithm for discovering clusters (DBSCAN),'' KDD 1996.%
}

\newcommand{\AMLLabSources}{%
M\"uller and Guido, \textit{Introduction to Machine Learning with Python}, companion notebooks.;%
McKinney, \textit{Python for Data Analysis} (3e), O'Reilly, 2022.;%
Scikit-learn user guide, accessed 2026-08-04.;%
Matplotlib and Seaborn documentation, accessed 2026-08-04.%
}
%
      }{%
        % Source strings for Applied Machine Learning (CSAI2017P) — used by slides + notes.
% Prescribed texts and reference books are taken from the course syllabus.
% Support texts are the author-hosted open-access editions:
%   statlearning.com            (ISL with Python)
%   probml.github.io/pml-book   (Murphy, Probabilistic Machine Learning)
%   d2l.ai                      (Dive into Deep Learning)
%   mml-book.github.io          (Mathematics for Machine Learning)

\newcommand{\AMLRepoURL}{https://github.com/tali7c/Applied-Machine-Learning}

% --- prescribed -------------------------------------------------------------
\newcommand{\AMLPrescribed}{%
M\"uller and Guido, \textit{Introduction to Machine Learning with Python}, O'Reilly, 2016.;%
Bishop, \textit{Pattern Recognition and Machine Learning}, Springer, 2006.;%
Murphy, \textit{Machine Learning: A Probabilistic Perspective}, MIT Press, 2012.;%
Han, Kamber and Pei, \textit{Data Mining: Concepts and Techniques} (3e), Morgan Kaufmann, 2012.%
}

% --- unit-wise ---------------------------------------------------------------
\newcommand{\AMLUnitISources}{%
M\"uller and Guido, \textit{Introduction to Machine Learning with Python}, Ch. 1.;%
Bishop, \textit{Pattern Recognition and Machine Learning}, Ch. 1.;%
Murphy, \textit{Probabilistic Machine Learning: An Introduction}, MIT Press, 2022, Ch. 1.;%
Mitchell, \textit{Machine Learning}, McGraw Hill, 1997, Ch. 1.;%
Scikit-learn documentation, accessed 2026-08-04.%
}

\newcommand{\AMLUnitIISources}{%
Bishop, \textit{Pattern Recognition and Machine Learning}, \S1.5, \S4.3, \S7.1.;%
Zhang, Lipton, Li and Smola, \textit{Dive into Deep Learning}, \S3.4.;%
Shalev-Shwartz and Ben-David, \textit{Understanding Machine Learning}, Ch. 15.;%
Murphy, \textit{Probabilistic Machine Learning: An Introduction}, Ch. 5.%
}

\newcommand{\AMLUnitIIISources}{%
Zhang, Lipton, Li and Smola, \textit{Dive into Deep Learning}, Ch. 12 (Optimization Algorithms).;%
Deisenroth, Faisal and Ong, \textit{Mathematics for Machine Learning}, Ch. 7.;%
Kingma and Ba, ``Adam: A Method for Stochastic Optimization,'' ICLR 2015.;%
Ruder, ``An overview of gradient descent optimization algorithms,'' arXiv:1609.04747, 2016.%
}

\newcommand{\AMLUnitIVSources}{%
Han, Kamber and Pei, \textit{Data Mining: Concepts and Techniques} (3e), Ch. 3.;%
M\"uller and Guido, \textit{Introduction to Machine Learning with Python}, Ch. 3--4.;%
James, Witten, Hastie, Tibshirani and Taylor, \textit{An Introduction to Statistical Learning with Python}, Ch. 5--6, 12.;%
Scikit-learn documentation (preprocessing, model selection), accessed 2026-08-04.%
}

\newcommand{\AMLUnitVSources}{%
James et al., \textit{An Introduction to Statistical Learning with Python}, Ch. 3--6.;%
Bishop, \textit{Pattern Recognition and Machine Learning}, Ch. 3.;%
Deisenroth, Faisal and Ong, \textit{Mathematics for Machine Learning}, Ch. 9.;%
M\"uller and Guido, \textit{Introduction to Machine Learning with Python}, \S2.3.%
}

\newcommand{\AMLUnitVISources}{%
Han, Kamber and Pei, \textit{Data Mining: Concepts and Techniques} (3e), Ch. 8--9.;%
James et al., \textit{An Introduction to Statistical Learning with Python}, Ch. 8--9.;%
Bishop, \textit{Pattern Recognition and Machine Learning}, Ch. 4--5, 7, 14.;%
Zhang et al., \textit{Dive into Deep Learning}, Ch. 5.%
}

\newcommand{\AMLUnitVIISources}{%
Han, Kamber and Pei, \textit{Data Mining: Concepts and Techniques} (3e), Ch. 10.;%
James et al., \textit{An Introduction to Statistical Learning with Python}, Ch. 12.;%
M\"uller and Guido, \textit{Introduction to Machine Learning with Python}, \S3.5.;%
Ester, Kriegel, Sander and Xu, ``A density-based algorithm for discovering clusters (DBSCAN),'' KDD 1996.%
}

\newcommand{\AMLLabSources}{%
M\"uller and Guido, \textit{Introduction to Machine Learning with Python}, companion notebooks.;%
McKinney, \textit{Python for Data Analysis} (3e), O'Reilly, 2022.;%
Scikit-learn user guide, accessed 2026-08-04.;%
Matplotlib and Seaborn documentation, accessed 2026-08-04.%
}
%
      }%
    }%
  }%
}
%
}{%
  \IfFileExists{../../../../../public/_shared/notes-common.tex}{%
    % Shared notes settings for Applied Machine Learning lectures (UPES).
% Keep this file free of \documentclass to allow reuse via \input.

\usepackage[utf8]{inputenc}
\usepackage[T1]{fontenc}
\usepackage{lmodern}          % scalable T1 fonts; without this pdflatex falls back
\input{glyphtounicode}        % to bitmapped EC Type 3 and the PDFs are neither
\pdfgentounicode=1            % searchable nor screen-reader friendly
\usepackage[a4paper,margin=1in]{geometry}
\usepackage{hyperref}
\usepackage{xurl}
\usepackage{booktabs}
\usepackage{amsmath}
\usepackage{amssymb}
\usepackage{listings}
\usepackage{xcolor}
\usepackage{enumitem}
\usepackage{parskip}

% Code listings (Python-oriented for this subject).
\lstset{
  language=Python,
  basicstyle=\ttfamily\small,
  keywordstyle=\color{blue},
  commentstyle=\color{gray},
  stringstyle=\color{teal},
  showstringspaces=false,
  columns=fullflexible,
  frame=single,
  framerule=0.2pt,
  breaklines=true
}

\setlist[itemize]{topsep=0.3em,itemsep=0.2em}
\setlist[enumerate]{topsep=0.3em,itemsep=0.2em}

% Simple per-section source line.
\newcommand{\notesources}[1]{\par\noindent{\small\color{gray}\textit{Sources:} \sloppy #1\par}}

% Unit-wise source strings for this subject (used by slides + notes).
% Path is relative to the lecture `latex/` folder (the compilation CWD).
\IfFileExists{../../../_shared/aml-sources.tex}{%
  % Source strings for Applied Machine Learning (CSAI2017P) — used by slides + notes.
% Prescribed texts and reference books are taken from the course syllabus.
% Support texts are the author-hosted open-access editions:
%   statlearning.com            (ISL with Python)
%   probml.github.io/pml-book   (Murphy, Probabilistic Machine Learning)
%   d2l.ai                      (Dive into Deep Learning)
%   mml-book.github.io          (Mathematics for Machine Learning)

\newcommand{\AMLRepoURL}{https://github.com/tali7c/Applied-Machine-Learning}

% --- prescribed -------------------------------------------------------------
\newcommand{\AMLPrescribed}{%
M\"uller and Guido, \textit{Introduction to Machine Learning with Python}, O'Reilly, 2016.;%
Bishop, \textit{Pattern Recognition and Machine Learning}, Springer, 2006.;%
Murphy, \textit{Machine Learning: A Probabilistic Perspective}, MIT Press, 2012.;%
Han, Kamber and Pei, \textit{Data Mining: Concepts and Techniques} (3e), Morgan Kaufmann, 2012.%
}

% --- unit-wise ---------------------------------------------------------------
\newcommand{\AMLUnitISources}{%
M\"uller and Guido, \textit{Introduction to Machine Learning with Python}, Ch. 1.;%
Bishop, \textit{Pattern Recognition and Machine Learning}, Ch. 1.;%
Murphy, \textit{Probabilistic Machine Learning: An Introduction}, MIT Press, 2022, Ch. 1.;%
Mitchell, \textit{Machine Learning}, McGraw Hill, 1997, Ch. 1.;%
Scikit-learn documentation, accessed 2026-08-04.%
}

\newcommand{\AMLUnitIISources}{%
Bishop, \textit{Pattern Recognition and Machine Learning}, \S1.5, \S4.3, \S7.1.;%
Zhang, Lipton, Li and Smola, \textit{Dive into Deep Learning}, \S3.4.;%
Shalev-Shwartz and Ben-David, \textit{Understanding Machine Learning}, Ch. 15.;%
Murphy, \textit{Probabilistic Machine Learning: An Introduction}, Ch. 5.%
}

\newcommand{\AMLUnitIIISources}{%
Zhang, Lipton, Li and Smola, \textit{Dive into Deep Learning}, Ch. 12 (Optimization Algorithms).;%
Deisenroth, Faisal and Ong, \textit{Mathematics for Machine Learning}, Ch. 7.;%
Kingma and Ba, ``Adam: A Method for Stochastic Optimization,'' ICLR 2015.;%
Ruder, ``An overview of gradient descent optimization algorithms,'' arXiv:1609.04747, 2016.%
}

\newcommand{\AMLUnitIVSources}{%
Han, Kamber and Pei, \textit{Data Mining: Concepts and Techniques} (3e), Ch. 3.;%
M\"uller and Guido, \textit{Introduction to Machine Learning with Python}, Ch. 3--4.;%
James, Witten, Hastie, Tibshirani and Taylor, \textit{An Introduction to Statistical Learning with Python}, Ch. 5--6, 12.;%
Scikit-learn documentation (preprocessing, model selection), accessed 2026-08-04.%
}

\newcommand{\AMLUnitVSources}{%
James et al., \textit{An Introduction to Statistical Learning with Python}, Ch. 3--6.;%
Bishop, \textit{Pattern Recognition and Machine Learning}, Ch. 3.;%
Deisenroth, Faisal and Ong, \textit{Mathematics for Machine Learning}, Ch. 9.;%
M\"uller and Guido, \textit{Introduction to Machine Learning with Python}, \S2.3.%
}

\newcommand{\AMLUnitVISources}{%
Han, Kamber and Pei, \textit{Data Mining: Concepts and Techniques} (3e), Ch. 8--9.;%
James et al., \textit{An Introduction to Statistical Learning with Python}, Ch. 8--9.;%
Bishop, \textit{Pattern Recognition and Machine Learning}, Ch. 4--5, 7, 14.;%
Zhang et al., \textit{Dive into Deep Learning}, Ch. 5.%
}

\newcommand{\AMLUnitVIISources}{%
Han, Kamber and Pei, \textit{Data Mining: Concepts and Techniques} (3e), Ch. 10.;%
James et al., \textit{An Introduction to Statistical Learning with Python}, Ch. 12.;%
M\"uller and Guido, \textit{Introduction to Machine Learning with Python}, \S3.5.;%
Ester, Kriegel, Sander and Xu, ``A density-based algorithm for discovering clusters (DBSCAN),'' KDD 1996.%
}

\newcommand{\AMLLabSources}{%
M\"uller and Guido, \textit{Introduction to Machine Learning with Python}, companion notebooks.;%
McKinney, \textit{Python for Data Analysis} (3e), O'Reilly, 2022.;%
Scikit-learn user guide, accessed 2026-08-04.;%
Matplotlib and Seaborn documentation, accessed 2026-08-04.%
}
%
}{%
  % Faculty copies live under `private/` so their compile CWD is different.
  \IfFileExists{../../../../public/_shared/aml-sources.tex}{%
    % Source strings for Applied Machine Learning (CSAI2017P) — used by slides + notes.
% Prescribed texts and reference books are taken from the course syllabus.
% Support texts are the author-hosted open-access editions:
%   statlearning.com            (ISL with Python)
%   probml.github.io/pml-book   (Murphy, Probabilistic Machine Learning)
%   d2l.ai                      (Dive into Deep Learning)
%   mml-book.github.io          (Mathematics for Machine Learning)

\newcommand{\AMLRepoURL}{https://github.com/tali7c/Applied-Machine-Learning}

% --- prescribed -------------------------------------------------------------
\newcommand{\AMLPrescribed}{%
M\"uller and Guido, \textit{Introduction to Machine Learning with Python}, O'Reilly, 2016.;%
Bishop, \textit{Pattern Recognition and Machine Learning}, Springer, 2006.;%
Murphy, \textit{Machine Learning: A Probabilistic Perspective}, MIT Press, 2012.;%
Han, Kamber and Pei, \textit{Data Mining: Concepts and Techniques} (3e), Morgan Kaufmann, 2012.%
}

% --- unit-wise ---------------------------------------------------------------
\newcommand{\AMLUnitISources}{%
M\"uller and Guido, \textit{Introduction to Machine Learning with Python}, Ch. 1.;%
Bishop, \textit{Pattern Recognition and Machine Learning}, Ch. 1.;%
Murphy, \textit{Probabilistic Machine Learning: An Introduction}, MIT Press, 2022, Ch. 1.;%
Mitchell, \textit{Machine Learning}, McGraw Hill, 1997, Ch. 1.;%
Scikit-learn documentation, accessed 2026-08-04.%
}

\newcommand{\AMLUnitIISources}{%
Bishop, \textit{Pattern Recognition and Machine Learning}, \S1.5, \S4.3, \S7.1.;%
Zhang, Lipton, Li and Smola, \textit{Dive into Deep Learning}, \S3.4.;%
Shalev-Shwartz and Ben-David, \textit{Understanding Machine Learning}, Ch. 15.;%
Murphy, \textit{Probabilistic Machine Learning: An Introduction}, Ch. 5.%
}

\newcommand{\AMLUnitIIISources}{%
Zhang, Lipton, Li and Smola, \textit{Dive into Deep Learning}, Ch. 12 (Optimization Algorithms).;%
Deisenroth, Faisal and Ong, \textit{Mathematics for Machine Learning}, Ch. 7.;%
Kingma and Ba, ``Adam: A Method for Stochastic Optimization,'' ICLR 2015.;%
Ruder, ``An overview of gradient descent optimization algorithms,'' arXiv:1609.04747, 2016.%
}

\newcommand{\AMLUnitIVSources}{%
Han, Kamber and Pei, \textit{Data Mining: Concepts and Techniques} (3e), Ch. 3.;%
M\"uller and Guido, \textit{Introduction to Machine Learning with Python}, Ch. 3--4.;%
James, Witten, Hastie, Tibshirani and Taylor, \textit{An Introduction to Statistical Learning with Python}, Ch. 5--6, 12.;%
Scikit-learn documentation (preprocessing, model selection), accessed 2026-08-04.%
}

\newcommand{\AMLUnitVSources}{%
James et al., \textit{An Introduction to Statistical Learning with Python}, Ch. 3--6.;%
Bishop, \textit{Pattern Recognition and Machine Learning}, Ch. 3.;%
Deisenroth, Faisal and Ong, \textit{Mathematics for Machine Learning}, Ch. 9.;%
M\"uller and Guido, \textit{Introduction to Machine Learning with Python}, \S2.3.%
}

\newcommand{\AMLUnitVISources}{%
Han, Kamber and Pei, \textit{Data Mining: Concepts and Techniques} (3e), Ch. 8--9.;%
James et al., \textit{An Introduction to Statistical Learning with Python}, Ch. 8--9.;%
Bishop, \textit{Pattern Recognition and Machine Learning}, Ch. 4--5, 7, 14.;%
Zhang et al., \textit{Dive into Deep Learning}, Ch. 5.%
}

\newcommand{\AMLUnitVIISources}{%
Han, Kamber and Pei, \textit{Data Mining: Concepts and Techniques} (3e), Ch. 10.;%
James et al., \textit{An Introduction to Statistical Learning with Python}, Ch. 12.;%
M\"uller and Guido, \textit{Introduction to Machine Learning with Python}, \S3.5.;%
Ester, Kriegel, Sander and Xu, ``A density-based algorithm for discovering clusters (DBSCAN),'' KDD 1996.%
}

\newcommand{\AMLLabSources}{%
M\"uller and Guido, \textit{Introduction to Machine Learning with Python}, companion notebooks.;%
McKinney, \textit{Python for Data Analysis} (3e), O'Reilly, 2022.;%
Scikit-learn user guide, accessed 2026-08-04.;%
Matplotlib and Seaborn documentation, accessed 2026-08-04.%
}
%
  }{%
    % Question papers live at `private/<family>/<instrument>/`.
    \IfFileExists{../../../public/_shared/aml-sources.tex}{%
      % Source strings for Applied Machine Learning (CSAI2017P) — used by slides + notes.
% Prescribed texts and reference books are taken from the course syllabus.
% Support texts are the author-hosted open-access editions:
%   statlearning.com            (ISL with Python)
%   probml.github.io/pml-book   (Murphy, Probabilistic Machine Learning)
%   d2l.ai                      (Dive into Deep Learning)
%   mml-book.github.io          (Mathematics for Machine Learning)

\newcommand{\AMLRepoURL}{https://github.com/tali7c/Applied-Machine-Learning}

% --- prescribed -------------------------------------------------------------
\newcommand{\AMLPrescribed}{%
M\"uller and Guido, \textit{Introduction to Machine Learning with Python}, O'Reilly, 2016.;%
Bishop, \textit{Pattern Recognition and Machine Learning}, Springer, 2006.;%
Murphy, \textit{Machine Learning: A Probabilistic Perspective}, MIT Press, 2012.;%
Han, Kamber and Pei, \textit{Data Mining: Concepts and Techniques} (3e), Morgan Kaufmann, 2012.%
}

% --- unit-wise ---------------------------------------------------------------
\newcommand{\AMLUnitISources}{%
M\"uller and Guido, \textit{Introduction to Machine Learning with Python}, Ch. 1.;%
Bishop, \textit{Pattern Recognition and Machine Learning}, Ch. 1.;%
Murphy, \textit{Probabilistic Machine Learning: An Introduction}, MIT Press, 2022, Ch. 1.;%
Mitchell, \textit{Machine Learning}, McGraw Hill, 1997, Ch. 1.;%
Scikit-learn documentation, accessed 2026-08-04.%
}

\newcommand{\AMLUnitIISources}{%
Bishop, \textit{Pattern Recognition and Machine Learning}, \S1.5, \S4.3, \S7.1.;%
Zhang, Lipton, Li and Smola, \textit{Dive into Deep Learning}, \S3.4.;%
Shalev-Shwartz and Ben-David, \textit{Understanding Machine Learning}, Ch. 15.;%
Murphy, \textit{Probabilistic Machine Learning: An Introduction}, Ch. 5.%
}

\newcommand{\AMLUnitIIISources}{%
Zhang, Lipton, Li and Smola, \textit{Dive into Deep Learning}, Ch. 12 (Optimization Algorithms).;%
Deisenroth, Faisal and Ong, \textit{Mathematics for Machine Learning}, Ch. 7.;%
Kingma and Ba, ``Adam: A Method for Stochastic Optimization,'' ICLR 2015.;%
Ruder, ``An overview of gradient descent optimization algorithms,'' arXiv:1609.04747, 2016.%
}

\newcommand{\AMLUnitIVSources}{%
Han, Kamber and Pei, \textit{Data Mining: Concepts and Techniques} (3e), Ch. 3.;%
M\"uller and Guido, \textit{Introduction to Machine Learning with Python}, Ch. 3--4.;%
James, Witten, Hastie, Tibshirani and Taylor, \textit{An Introduction to Statistical Learning with Python}, Ch. 5--6, 12.;%
Scikit-learn documentation (preprocessing, model selection), accessed 2026-08-04.%
}

\newcommand{\AMLUnitVSources}{%
James et al., \textit{An Introduction to Statistical Learning with Python}, Ch. 3--6.;%
Bishop, \textit{Pattern Recognition and Machine Learning}, Ch. 3.;%
Deisenroth, Faisal and Ong, \textit{Mathematics for Machine Learning}, Ch. 9.;%
M\"uller and Guido, \textit{Introduction to Machine Learning with Python}, \S2.3.%
}

\newcommand{\AMLUnitVISources}{%
Han, Kamber and Pei, \textit{Data Mining: Concepts and Techniques} (3e), Ch. 8--9.;%
James et al., \textit{An Introduction to Statistical Learning with Python}, Ch. 8--9.;%
Bishop, \textit{Pattern Recognition and Machine Learning}, Ch. 4--5, 7, 14.;%
Zhang et al., \textit{Dive into Deep Learning}, Ch. 5.%
}

\newcommand{\AMLUnitVIISources}{%
Han, Kamber and Pei, \textit{Data Mining: Concepts and Techniques} (3e), Ch. 10.;%
James et al., \textit{An Introduction to Statistical Learning with Python}, Ch. 12.;%
M\"uller and Guido, \textit{Introduction to Machine Learning with Python}, \S3.5.;%
Ester, Kriegel, Sander and Xu, ``A density-based algorithm for discovering clusters (DBSCAN),'' KDD 1996.%
}

\newcommand{\AMLLabSources}{%
M\"uller and Guido, \textit{Introduction to Machine Learning with Python}, companion notebooks.;%
McKinney, \textit{Python for Data Analysis} (3e), O'Reilly, 2022.;%
Scikit-learn user guide, accessed 2026-08-04.;%
Matplotlib and Seaborn documentation, accessed 2026-08-04.%
}
%
    }{%
      % Marking schemes live at `private/solutions/<family>/<id>/latex/`.
      \IfFileExists{../../../../../public/_shared/aml-sources.tex}{%
        % Source strings for Applied Machine Learning (CSAI2017P) — used by slides + notes.
% Prescribed texts and reference books are taken from the course syllabus.
% Support texts are the author-hosted open-access editions:
%   statlearning.com            (ISL with Python)
%   probml.github.io/pml-book   (Murphy, Probabilistic Machine Learning)
%   d2l.ai                      (Dive into Deep Learning)
%   mml-book.github.io          (Mathematics for Machine Learning)

\newcommand{\AMLRepoURL}{https://github.com/tali7c/Applied-Machine-Learning}

% --- prescribed -------------------------------------------------------------
\newcommand{\AMLPrescribed}{%
M\"uller and Guido, \textit{Introduction to Machine Learning with Python}, O'Reilly, 2016.;%
Bishop, \textit{Pattern Recognition and Machine Learning}, Springer, 2006.;%
Murphy, \textit{Machine Learning: A Probabilistic Perspective}, MIT Press, 2012.;%
Han, Kamber and Pei, \textit{Data Mining: Concepts and Techniques} (3e), Morgan Kaufmann, 2012.%
}

% --- unit-wise ---------------------------------------------------------------
\newcommand{\AMLUnitISources}{%
M\"uller and Guido, \textit{Introduction to Machine Learning with Python}, Ch. 1.;%
Bishop, \textit{Pattern Recognition and Machine Learning}, Ch. 1.;%
Murphy, \textit{Probabilistic Machine Learning: An Introduction}, MIT Press, 2022, Ch. 1.;%
Mitchell, \textit{Machine Learning}, McGraw Hill, 1997, Ch. 1.;%
Scikit-learn documentation, accessed 2026-08-04.%
}

\newcommand{\AMLUnitIISources}{%
Bishop, \textit{Pattern Recognition and Machine Learning}, \S1.5, \S4.3, \S7.1.;%
Zhang, Lipton, Li and Smola, \textit{Dive into Deep Learning}, \S3.4.;%
Shalev-Shwartz and Ben-David, \textit{Understanding Machine Learning}, Ch. 15.;%
Murphy, \textit{Probabilistic Machine Learning: An Introduction}, Ch. 5.%
}

\newcommand{\AMLUnitIIISources}{%
Zhang, Lipton, Li and Smola, \textit{Dive into Deep Learning}, Ch. 12 (Optimization Algorithms).;%
Deisenroth, Faisal and Ong, \textit{Mathematics for Machine Learning}, Ch. 7.;%
Kingma and Ba, ``Adam: A Method for Stochastic Optimization,'' ICLR 2015.;%
Ruder, ``An overview of gradient descent optimization algorithms,'' arXiv:1609.04747, 2016.%
}

\newcommand{\AMLUnitIVSources}{%
Han, Kamber and Pei, \textit{Data Mining: Concepts and Techniques} (3e), Ch. 3.;%
M\"uller and Guido, \textit{Introduction to Machine Learning with Python}, Ch. 3--4.;%
James, Witten, Hastie, Tibshirani and Taylor, \textit{An Introduction to Statistical Learning with Python}, Ch. 5--6, 12.;%
Scikit-learn documentation (preprocessing, model selection), accessed 2026-08-04.%
}

\newcommand{\AMLUnitVSources}{%
James et al., \textit{An Introduction to Statistical Learning with Python}, Ch. 3--6.;%
Bishop, \textit{Pattern Recognition and Machine Learning}, Ch. 3.;%
Deisenroth, Faisal and Ong, \textit{Mathematics for Machine Learning}, Ch. 9.;%
M\"uller and Guido, \textit{Introduction to Machine Learning with Python}, \S2.3.%
}

\newcommand{\AMLUnitVISources}{%
Han, Kamber and Pei, \textit{Data Mining: Concepts and Techniques} (3e), Ch. 8--9.;%
James et al., \textit{An Introduction to Statistical Learning with Python}, Ch. 8--9.;%
Bishop, \textit{Pattern Recognition and Machine Learning}, Ch. 4--5, 7, 14.;%
Zhang et al., \textit{Dive into Deep Learning}, Ch. 5.%
}

\newcommand{\AMLUnitVIISources}{%
Han, Kamber and Pei, \textit{Data Mining: Concepts and Techniques} (3e), Ch. 10.;%
James et al., \textit{An Introduction to Statistical Learning with Python}, Ch. 12.;%
M\"uller and Guido, \textit{Introduction to Machine Learning with Python}, \S3.5.;%
Ester, Kriegel, Sander and Xu, ``A density-based algorithm for discovering clusters (DBSCAN),'' KDD 1996.%
}

\newcommand{\AMLLabSources}{%
M\"uller and Guido, \textit{Introduction to Machine Learning with Python}, companion notebooks.;%
McKinney, \textit{Python for Data Analysis} (3e), O'Reilly, 2022.;%
Scikit-learn user guide, accessed 2026-08-04.;%
Matplotlib and Seaborn documentation, accessed 2026-08-04.%
}
%
      }{%
        % Source strings for Applied Machine Learning (CSAI2017P) — used by slides + notes.
% Prescribed texts and reference books are taken from the course syllabus.
% Support texts are the author-hosted open-access editions:
%   statlearning.com            (ISL with Python)
%   probml.github.io/pml-book   (Murphy, Probabilistic Machine Learning)
%   d2l.ai                      (Dive into Deep Learning)
%   mml-book.github.io          (Mathematics for Machine Learning)

\newcommand{\AMLRepoURL}{https://github.com/tali7c/Applied-Machine-Learning}

% --- prescribed -------------------------------------------------------------
\newcommand{\AMLPrescribed}{%
M\"uller and Guido, \textit{Introduction to Machine Learning with Python}, O'Reilly, 2016.;%
Bishop, \textit{Pattern Recognition and Machine Learning}, Springer, 2006.;%
Murphy, \textit{Machine Learning: A Probabilistic Perspective}, MIT Press, 2012.;%
Han, Kamber and Pei, \textit{Data Mining: Concepts and Techniques} (3e), Morgan Kaufmann, 2012.%
}

% --- unit-wise ---------------------------------------------------------------
\newcommand{\AMLUnitISources}{%
M\"uller and Guido, \textit{Introduction to Machine Learning with Python}, Ch. 1.;%
Bishop, \textit{Pattern Recognition and Machine Learning}, Ch. 1.;%
Murphy, \textit{Probabilistic Machine Learning: An Introduction}, MIT Press, 2022, Ch. 1.;%
Mitchell, \textit{Machine Learning}, McGraw Hill, 1997, Ch. 1.;%
Scikit-learn documentation, accessed 2026-08-04.%
}

\newcommand{\AMLUnitIISources}{%
Bishop, \textit{Pattern Recognition and Machine Learning}, \S1.5, \S4.3, \S7.1.;%
Zhang, Lipton, Li and Smola, \textit{Dive into Deep Learning}, \S3.4.;%
Shalev-Shwartz and Ben-David, \textit{Understanding Machine Learning}, Ch. 15.;%
Murphy, \textit{Probabilistic Machine Learning: An Introduction}, Ch. 5.%
}

\newcommand{\AMLUnitIIISources}{%
Zhang, Lipton, Li and Smola, \textit{Dive into Deep Learning}, Ch. 12 (Optimization Algorithms).;%
Deisenroth, Faisal and Ong, \textit{Mathematics for Machine Learning}, Ch. 7.;%
Kingma and Ba, ``Adam: A Method for Stochastic Optimization,'' ICLR 2015.;%
Ruder, ``An overview of gradient descent optimization algorithms,'' arXiv:1609.04747, 2016.%
}

\newcommand{\AMLUnitIVSources}{%
Han, Kamber and Pei, \textit{Data Mining: Concepts and Techniques} (3e), Ch. 3.;%
M\"uller and Guido, \textit{Introduction to Machine Learning with Python}, Ch. 3--4.;%
James, Witten, Hastie, Tibshirani and Taylor, \textit{An Introduction to Statistical Learning with Python}, Ch. 5--6, 12.;%
Scikit-learn documentation (preprocessing, model selection), accessed 2026-08-04.%
}

\newcommand{\AMLUnitVSources}{%
James et al., \textit{An Introduction to Statistical Learning with Python}, Ch. 3--6.;%
Bishop, \textit{Pattern Recognition and Machine Learning}, Ch. 3.;%
Deisenroth, Faisal and Ong, \textit{Mathematics for Machine Learning}, Ch. 9.;%
M\"uller and Guido, \textit{Introduction to Machine Learning with Python}, \S2.3.%
}

\newcommand{\AMLUnitVISources}{%
Han, Kamber and Pei, \textit{Data Mining: Concepts and Techniques} (3e), Ch. 8--9.;%
James et al., \textit{An Introduction to Statistical Learning with Python}, Ch. 8--9.;%
Bishop, \textit{Pattern Recognition and Machine Learning}, Ch. 4--5, 7, 14.;%
Zhang et al., \textit{Dive into Deep Learning}, Ch. 5.%
}

\newcommand{\AMLUnitVIISources}{%
Han, Kamber and Pei, \textit{Data Mining: Concepts and Techniques} (3e), Ch. 10.;%
James et al., \textit{An Introduction to Statistical Learning with Python}, Ch. 12.;%
M\"uller and Guido, \textit{Introduction to Machine Learning with Python}, \S3.5.;%
Ester, Kriegel, Sander and Xu, ``A density-based algorithm for discovering clusters (DBSCAN),'' KDD 1996.%
}

\newcommand{\AMLLabSources}{%
M\"uller and Guido, \textit{Introduction to Machine Learning with Python}, companion notebooks.;%
McKinney, \textit{Python for Data Analysis} (3e), O'Reilly, 2022.;%
Scikit-learn user guide, accessed 2026-08-04.;%
Matplotlib and Seaborn documentation, accessed 2026-08-04.%
}
%
      }%
    }%
  }%
}
%
  }{%
    % Shared notes settings for Applied Machine Learning lectures (UPES).
% Keep this file free of \documentclass to allow reuse via \input.

\usepackage[utf8]{inputenc}
\usepackage[T1]{fontenc}
\usepackage{lmodern}          % scalable T1 fonts; without this pdflatex falls back
\input{glyphtounicode}        % to bitmapped EC Type 3 and the PDFs are neither
\pdfgentounicode=1            % searchable nor screen-reader friendly
\usepackage[a4paper,margin=1in]{geometry}
\usepackage{hyperref}
\usepackage{xurl}
\usepackage{booktabs}
\usepackage{amsmath}
\usepackage{amssymb}
\usepackage{listings}
\usepackage{xcolor}
\usepackage{enumitem}
\usepackage{parskip}

% Code listings (Python-oriented for this subject).
\lstset{
  language=Python,
  basicstyle=\ttfamily\small,
  keywordstyle=\color{blue},
  commentstyle=\color{gray},
  stringstyle=\color{teal},
  showstringspaces=false,
  columns=fullflexible,
  frame=single,
  framerule=0.2pt,
  breaklines=true
}

\setlist[itemize]{topsep=0.3em,itemsep=0.2em}
\setlist[enumerate]{topsep=0.3em,itemsep=0.2em}

% Simple per-section source line.
\newcommand{\notesources}[1]{\par\noindent{\small\color{gray}\textit{Sources:} \sloppy #1\par}}

% Unit-wise source strings for this subject (used by slides + notes).
% Path is relative to the lecture `latex/` folder (the compilation CWD).
\IfFileExists{../../../_shared/aml-sources.tex}{%
  % Source strings for Applied Machine Learning (CSAI2017P) — used by slides + notes.
% Prescribed texts and reference books are taken from the course syllabus.
% Support texts are the author-hosted open-access editions:
%   statlearning.com            (ISL with Python)
%   probml.github.io/pml-book   (Murphy, Probabilistic Machine Learning)
%   d2l.ai                      (Dive into Deep Learning)
%   mml-book.github.io          (Mathematics for Machine Learning)

\newcommand{\AMLRepoURL}{https://github.com/tali7c/Applied-Machine-Learning}

% --- prescribed -------------------------------------------------------------
\newcommand{\AMLPrescribed}{%
M\"uller and Guido, \textit{Introduction to Machine Learning with Python}, O'Reilly, 2016.;%
Bishop, \textit{Pattern Recognition and Machine Learning}, Springer, 2006.;%
Murphy, \textit{Machine Learning: A Probabilistic Perspective}, MIT Press, 2012.;%
Han, Kamber and Pei, \textit{Data Mining: Concepts and Techniques} (3e), Morgan Kaufmann, 2012.%
}

% --- unit-wise ---------------------------------------------------------------
\newcommand{\AMLUnitISources}{%
M\"uller and Guido, \textit{Introduction to Machine Learning with Python}, Ch. 1.;%
Bishop, \textit{Pattern Recognition and Machine Learning}, Ch. 1.;%
Murphy, \textit{Probabilistic Machine Learning: An Introduction}, MIT Press, 2022, Ch. 1.;%
Mitchell, \textit{Machine Learning}, McGraw Hill, 1997, Ch. 1.;%
Scikit-learn documentation, accessed 2026-08-04.%
}

\newcommand{\AMLUnitIISources}{%
Bishop, \textit{Pattern Recognition and Machine Learning}, \S1.5, \S4.3, \S7.1.;%
Zhang, Lipton, Li and Smola, \textit{Dive into Deep Learning}, \S3.4.;%
Shalev-Shwartz and Ben-David, \textit{Understanding Machine Learning}, Ch. 15.;%
Murphy, \textit{Probabilistic Machine Learning: An Introduction}, Ch. 5.%
}

\newcommand{\AMLUnitIIISources}{%
Zhang, Lipton, Li and Smola, \textit{Dive into Deep Learning}, Ch. 12 (Optimization Algorithms).;%
Deisenroth, Faisal and Ong, \textit{Mathematics for Machine Learning}, Ch. 7.;%
Kingma and Ba, ``Adam: A Method for Stochastic Optimization,'' ICLR 2015.;%
Ruder, ``An overview of gradient descent optimization algorithms,'' arXiv:1609.04747, 2016.%
}

\newcommand{\AMLUnitIVSources}{%
Han, Kamber and Pei, \textit{Data Mining: Concepts and Techniques} (3e), Ch. 3.;%
M\"uller and Guido, \textit{Introduction to Machine Learning with Python}, Ch. 3--4.;%
James, Witten, Hastie, Tibshirani and Taylor, \textit{An Introduction to Statistical Learning with Python}, Ch. 5--6, 12.;%
Scikit-learn documentation (preprocessing, model selection), accessed 2026-08-04.%
}

\newcommand{\AMLUnitVSources}{%
James et al., \textit{An Introduction to Statistical Learning with Python}, Ch. 3--6.;%
Bishop, \textit{Pattern Recognition and Machine Learning}, Ch. 3.;%
Deisenroth, Faisal and Ong, \textit{Mathematics for Machine Learning}, Ch. 9.;%
M\"uller and Guido, \textit{Introduction to Machine Learning with Python}, \S2.3.%
}

\newcommand{\AMLUnitVISources}{%
Han, Kamber and Pei, \textit{Data Mining: Concepts and Techniques} (3e), Ch. 8--9.;%
James et al., \textit{An Introduction to Statistical Learning with Python}, Ch. 8--9.;%
Bishop, \textit{Pattern Recognition and Machine Learning}, Ch. 4--5, 7, 14.;%
Zhang et al., \textit{Dive into Deep Learning}, Ch. 5.%
}

\newcommand{\AMLUnitVIISources}{%
Han, Kamber and Pei, \textit{Data Mining: Concepts and Techniques} (3e), Ch. 10.;%
James et al., \textit{An Introduction to Statistical Learning with Python}, Ch. 12.;%
M\"uller and Guido, \textit{Introduction to Machine Learning with Python}, \S3.5.;%
Ester, Kriegel, Sander and Xu, ``A density-based algorithm for discovering clusters (DBSCAN),'' KDD 1996.%
}

\newcommand{\AMLLabSources}{%
M\"uller and Guido, \textit{Introduction to Machine Learning with Python}, companion notebooks.;%
McKinney, \textit{Python for Data Analysis} (3e), O'Reilly, 2022.;%
Scikit-learn user guide, accessed 2026-08-04.;%
Matplotlib and Seaborn documentation, accessed 2026-08-04.%
}
%
}{%
  % Faculty copies live under `private/` so their compile CWD is different.
  \IfFileExists{../../../../public/_shared/aml-sources.tex}{%
    % Source strings for Applied Machine Learning (CSAI2017P) — used by slides + notes.
% Prescribed texts and reference books are taken from the course syllabus.
% Support texts are the author-hosted open-access editions:
%   statlearning.com            (ISL with Python)
%   probml.github.io/pml-book   (Murphy, Probabilistic Machine Learning)
%   d2l.ai                      (Dive into Deep Learning)
%   mml-book.github.io          (Mathematics for Machine Learning)

\newcommand{\AMLRepoURL}{https://github.com/tali7c/Applied-Machine-Learning}

% --- prescribed -------------------------------------------------------------
\newcommand{\AMLPrescribed}{%
M\"uller and Guido, \textit{Introduction to Machine Learning with Python}, O'Reilly, 2016.;%
Bishop, \textit{Pattern Recognition and Machine Learning}, Springer, 2006.;%
Murphy, \textit{Machine Learning: A Probabilistic Perspective}, MIT Press, 2012.;%
Han, Kamber and Pei, \textit{Data Mining: Concepts and Techniques} (3e), Morgan Kaufmann, 2012.%
}

% --- unit-wise ---------------------------------------------------------------
\newcommand{\AMLUnitISources}{%
M\"uller and Guido, \textit{Introduction to Machine Learning with Python}, Ch. 1.;%
Bishop, \textit{Pattern Recognition and Machine Learning}, Ch. 1.;%
Murphy, \textit{Probabilistic Machine Learning: An Introduction}, MIT Press, 2022, Ch. 1.;%
Mitchell, \textit{Machine Learning}, McGraw Hill, 1997, Ch. 1.;%
Scikit-learn documentation, accessed 2026-08-04.%
}

\newcommand{\AMLUnitIISources}{%
Bishop, \textit{Pattern Recognition and Machine Learning}, \S1.5, \S4.3, \S7.1.;%
Zhang, Lipton, Li and Smola, \textit{Dive into Deep Learning}, \S3.4.;%
Shalev-Shwartz and Ben-David, \textit{Understanding Machine Learning}, Ch. 15.;%
Murphy, \textit{Probabilistic Machine Learning: An Introduction}, Ch. 5.%
}

\newcommand{\AMLUnitIIISources}{%
Zhang, Lipton, Li and Smola, \textit{Dive into Deep Learning}, Ch. 12 (Optimization Algorithms).;%
Deisenroth, Faisal and Ong, \textit{Mathematics for Machine Learning}, Ch. 7.;%
Kingma and Ba, ``Adam: A Method for Stochastic Optimization,'' ICLR 2015.;%
Ruder, ``An overview of gradient descent optimization algorithms,'' arXiv:1609.04747, 2016.%
}

\newcommand{\AMLUnitIVSources}{%
Han, Kamber and Pei, \textit{Data Mining: Concepts and Techniques} (3e), Ch. 3.;%
M\"uller and Guido, \textit{Introduction to Machine Learning with Python}, Ch. 3--4.;%
James, Witten, Hastie, Tibshirani and Taylor, \textit{An Introduction to Statistical Learning with Python}, Ch. 5--6, 12.;%
Scikit-learn documentation (preprocessing, model selection), accessed 2026-08-04.%
}

\newcommand{\AMLUnitVSources}{%
James et al., \textit{An Introduction to Statistical Learning with Python}, Ch. 3--6.;%
Bishop, \textit{Pattern Recognition and Machine Learning}, Ch. 3.;%
Deisenroth, Faisal and Ong, \textit{Mathematics for Machine Learning}, Ch. 9.;%
M\"uller and Guido, \textit{Introduction to Machine Learning with Python}, \S2.3.%
}

\newcommand{\AMLUnitVISources}{%
Han, Kamber and Pei, \textit{Data Mining: Concepts and Techniques} (3e), Ch. 8--9.;%
James et al., \textit{An Introduction to Statistical Learning with Python}, Ch. 8--9.;%
Bishop, \textit{Pattern Recognition and Machine Learning}, Ch. 4--5, 7, 14.;%
Zhang et al., \textit{Dive into Deep Learning}, Ch. 5.%
}

\newcommand{\AMLUnitVIISources}{%
Han, Kamber and Pei, \textit{Data Mining: Concepts and Techniques} (3e), Ch. 10.;%
James et al., \textit{An Introduction to Statistical Learning with Python}, Ch. 12.;%
M\"uller and Guido, \textit{Introduction to Machine Learning with Python}, \S3.5.;%
Ester, Kriegel, Sander and Xu, ``A density-based algorithm for discovering clusters (DBSCAN),'' KDD 1996.%
}

\newcommand{\AMLLabSources}{%
M\"uller and Guido, \textit{Introduction to Machine Learning with Python}, companion notebooks.;%
McKinney, \textit{Python for Data Analysis} (3e), O'Reilly, 2022.;%
Scikit-learn user guide, accessed 2026-08-04.;%
Matplotlib and Seaborn documentation, accessed 2026-08-04.%
}
%
  }{%
    % Question papers live at `private/<family>/<instrument>/`.
    \IfFileExists{../../../public/_shared/aml-sources.tex}{%
      % Source strings for Applied Machine Learning (CSAI2017P) — used by slides + notes.
% Prescribed texts and reference books are taken from the course syllabus.
% Support texts are the author-hosted open-access editions:
%   statlearning.com            (ISL with Python)
%   probml.github.io/pml-book   (Murphy, Probabilistic Machine Learning)
%   d2l.ai                      (Dive into Deep Learning)
%   mml-book.github.io          (Mathematics for Machine Learning)

\newcommand{\AMLRepoURL}{https://github.com/tali7c/Applied-Machine-Learning}

% --- prescribed -------------------------------------------------------------
\newcommand{\AMLPrescribed}{%
M\"uller and Guido, \textit{Introduction to Machine Learning with Python}, O'Reilly, 2016.;%
Bishop, \textit{Pattern Recognition and Machine Learning}, Springer, 2006.;%
Murphy, \textit{Machine Learning: A Probabilistic Perspective}, MIT Press, 2012.;%
Han, Kamber and Pei, \textit{Data Mining: Concepts and Techniques} (3e), Morgan Kaufmann, 2012.%
}

% --- unit-wise ---------------------------------------------------------------
\newcommand{\AMLUnitISources}{%
M\"uller and Guido, \textit{Introduction to Machine Learning with Python}, Ch. 1.;%
Bishop, \textit{Pattern Recognition and Machine Learning}, Ch. 1.;%
Murphy, \textit{Probabilistic Machine Learning: An Introduction}, MIT Press, 2022, Ch. 1.;%
Mitchell, \textit{Machine Learning}, McGraw Hill, 1997, Ch. 1.;%
Scikit-learn documentation, accessed 2026-08-04.%
}

\newcommand{\AMLUnitIISources}{%
Bishop, \textit{Pattern Recognition and Machine Learning}, \S1.5, \S4.3, \S7.1.;%
Zhang, Lipton, Li and Smola, \textit{Dive into Deep Learning}, \S3.4.;%
Shalev-Shwartz and Ben-David, \textit{Understanding Machine Learning}, Ch. 15.;%
Murphy, \textit{Probabilistic Machine Learning: An Introduction}, Ch. 5.%
}

\newcommand{\AMLUnitIIISources}{%
Zhang, Lipton, Li and Smola, \textit{Dive into Deep Learning}, Ch. 12 (Optimization Algorithms).;%
Deisenroth, Faisal and Ong, \textit{Mathematics for Machine Learning}, Ch. 7.;%
Kingma and Ba, ``Adam: A Method for Stochastic Optimization,'' ICLR 2015.;%
Ruder, ``An overview of gradient descent optimization algorithms,'' arXiv:1609.04747, 2016.%
}

\newcommand{\AMLUnitIVSources}{%
Han, Kamber and Pei, \textit{Data Mining: Concepts and Techniques} (3e), Ch. 3.;%
M\"uller and Guido, \textit{Introduction to Machine Learning with Python}, Ch. 3--4.;%
James, Witten, Hastie, Tibshirani and Taylor, \textit{An Introduction to Statistical Learning with Python}, Ch. 5--6, 12.;%
Scikit-learn documentation (preprocessing, model selection), accessed 2026-08-04.%
}

\newcommand{\AMLUnitVSources}{%
James et al., \textit{An Introduction to Statistical Learning with Python}, Ch. 3--6.;%
Bishop, \textit{Pattern Recognition and Machine Learning}, Ch. 3.;%
Deisenroth, Faisal and Ong, \textit{Mathematics for Machine Learning}, Ch. 9.;%
M\"uller and Guido, \textit{Introduction to Machine Learning with Python}, \S2.3.%
}

\newcommand{\AMLUnitVISources}{%
Han, Kamber and Pei, \textit{Data Mining: Concepts and Techniques} (3e), Ch. 8--9.;%
James et al., \textit{An Introduction to Statistical Learning with Python}, Ch. 8--9.;%
Bishop, \textit{Pattern Recognition and Machine Learning}, Ch. 4--5, 7, 14.;%
Zhang et al., \textit{Dive into Deep Learning}, Ch. 5.%
}

\newcommand{\AMLUnitVIISources}{%
Han, Kamber and Pei, \textit{Data Mining: Concepts and Techniques} (3e), Ch. 10.;%
James et al., \textit{An Introduction to Statistical Learning with Python}, Ch. 12.;%
M\"uller and Guido, \textit{Introduction to Machine Learning with Python}, \S3.5.;%
Ester, Kriegel, Sander and Xu, ``A density-based algorithm for discovering clusters (DBSCAN),'' KDD 1996.%
}

\newcommand{\AMLLabSources}{%
M\"uller and Guido, \textit{Introduction to Machine Learning with Python}, companion notebooks.;%
McKinney, \textit{Python for Data Analysis} (3e), O'Reilly, 2022.;%
Scikit-learn user guide, accessed 2026-08-04.;%
Matplotlib and Seaborn documentation, accessed 2026-08-04.%
}
%
    }{%
      % Marking schemes live at `private/solutions/<family>/<id>/latex/`.
      \IfFileExists{../../../../../public/_shared/aml-sources.tex}{%
        % Source strings for Applied Machine Learning (CSAI2017P) — used by slides + notes.
% Prescribed texts and reference books are taken from the course syllabus.
% Support texts are the author-hosted open-access editions:
%   statlearning.com            (ISL with Python)
%   probml.github.io/pml-book   (Murphy, Probabilistic Machine Learning)
%   d2l.ai                      (Dive into Deep Learning)
%   mml-book.github.io          (Mathematics for Machine Learning)

\newcommand{\AMLRepoURL}{https://github.com/tali7c/Applied-Machine-Learning}

% --- prescribed -------------------------------------------------------------
\newcommand{\AMLPrescribed}{%
M\"uller and Guido, \textit{Introduction to Machine Learning with Python}, O'Reilly, 2016.;%
Bishop, \textit{Pattern Recognition and Machine Learning}, Springer, 2006.;%
Murphy, \textit{Machine Learning: A Probabilistic Perspective}, MIT Press, 2012.;%
Han, Kamber and Pei, \textit{Data Mining: Concepts and Techniques} (3e), Morgan Kaufmann, 2012.%
}

% --- unit-wise ---------------------------------------------------------------
\newcommand{\AMLUnitISources}{%
M\"uller and Guido, \textit{Introduction to Machine Learning with Python}, Ch. 1.;%
Bishop, \textit{Pattern Recognition and Machine Learning}, Ch. 1.;%
Murphy, \textit{Probabilistic Machine Learning: An Introduction}, MIT Press, 2022, Ch. 1.;%
Mitchell, \textit{Machine Learning}, McGraw Hill, 1997, Ch. 1.;%
Scikit-learn documentation, accessed 2026-08-04.%
}

\newcommand{\AMLUnitIISources}{%
Bishop, \textit{Pattern Recognition and Machine Learning}, \S1.5, \S4.3, \S7.1.;%
Zhang, Lipton, Li and Smola, \textit{Dive into Deep Learning}, \S3.4.;%
Shalev-Shwartz and Ben-David, \textit{Understanding Machine Learning}, Ch. 15.;%
Murphy, \textit{Probabilistic Machine Learning: An Introduction}, Ch. 5.%
}

\newcommand{\AMLUnitIIISources}{%
Zhang, Lipton, Li and Smola, \textit{Dive into Deep Learning}, Ch. 12 (Optimization Algorithms).;%
Deisenroth, Faisal and Ong, \textit{Mathematics for Machine Learning}, Ch. 7.;%
Kingma and Ba, ``Adam: A Method for Stochastic Optimization,'' ICLR 2015.;%
Ruder, ``An overview of gradient descent optimization algorithms,'' arXiv:1609.04747, 2016.%
}

\newcommand{\AMLUnitIVSources}{%
Han, Kamber and Pei, \textit{Data Mining: Concepts and Techniques} (3e), Ch. 3.;%
M\"uller and Guido, \textit{Introduction to Machine Learning with Python}, Ch. 3--4.;%
James, Witten, Hastie, Tibshirani and Taylor, \textit{An Introduction to Statistical Learning with Python}, Ch. 5--6, 12.;%
Scikit-learn documentation (preprocessing, model selection), accessed 2026-08-04.%
}

\newcommand{\AMLUnitVSources}{%
James et al., \textit{An Introduction to Statistical Learning with Python}, Ch. 3--6.;%
Bishop, \textit{Pattern Recognition and Machine Learning}, Ch. 3.;%
Deisenroth, Faisal and Ong, \textit{Mathematics for Machine Learning}, Ch. 9.;%
M\"uller and Guido, \textit{Introduction to Machine Learning with Python}, \S2.3.%
}

\newcommand{\AMLUnitVISources}{%
Han, Kamber and Pei, \textit{Data Mining: Concepts and Techniques} (3e), Ch. 8--9.;%
James et al., \textit{An Introduction to Statistical Learning with Python}, Ch. 8--9.;%
Bishop, \textit{Pattern Recognition and Machine Learning}, Ch. 4--5, 7, 14.;%
Zhang et al., \textit{Dive into Deep Learning}, Ch. 5.%
}

\newcommand{\AMLUnitVIISources}{%
Han, Kamber and Pei, \textit{Data Mining: Concepts and Techniques} (3e), Ch. 10.;%
James et al., \textit{An Introduction to Statistical Learning with Python}, Ch. 12.;%
M\"uller and Guido, \textit{Introduction to Machine Learning with Python}, \S3.5.;%
Ester, Kriegel, Sander and Xu, ``A density-based algorithm for discovering clusters (DBSCAN),'' KDD 1996.%
}

\newcommand{\AMLLabSources}{%
M\"uller and Guido, \textit{Introduction to Machine Learning with Python}, companion notebooks.;%
McKinney, \textit{Python for Data Analysis} (3e), O'Reilly, 2022.;%
Scikit-learn user guide, accessed 2026-08-04.;%
Matplotlib and Seaborn documentation, accessed 2026-08-04.%
}
%
      }{%
        % Source strings for Applied Machine Learning (CSAI2017P) — used by slides + notes.
% Prescribed texts and reference books are taken from the course syllabus.
% Support texts are the author-hosted open-access editions:
%   statlearning.com            (ISL with Python)
%   probml.github.io/pml-book   (Murphy, Probabilistic Machine Learning)
%   d2l.ai                      (Dive into Deep Learning)
%   mml-book.github.io          (Mathematics for Machine Learning)

\newcommand{\AMLRepoURL}{https://github.com/tali7c/Applied-Machine-Learning}

% --- prescribed -------------------------------------------------------------
\newcommand{\AMLPrescribed}{%
M\"uller and Guido, \textit{Introduction to Machine Learning with Python}, O'Reilly, 2016.;%
Bishop, \textit{Pattern Recognition and Machine Learning}, Springer, 2006.;%
Murphy, \textit{Machine Learning: A Probabilistic Perspective}, MIT Press, 2012.;%
Han, Kamber and Pei, \textit{Data Mining: Concepts and Techniques} (3e), Morgan Kaufmann, 2012.%
}

% --- unit-wise ---------------------------------------------------------------
\newcommand{\AMLUnitISources}{%
M\"uller and Guido, \textit{Introduction to Machine Learning with Python}, Ch. 1.;%
Bishop, \textit{Pattern Recognition and Machine Learning}, Ch. 1.;%
Murphy, \textit{Probabilistic Machine Learning: An Introduction}, MIT Press, 2022, Ch. 1.;%
Mitchell, \textit{Machine Learning}, McGraw Hill, 1997, Ch. 1.;%
Scikit-learn documentation, accessed 2026-08-04.%
}

\newcommand{\AMLUnitIISources}{%
Bishop, \textit{Pattern Recognition and Machine Learning}, \S1.5, \S4.3, \S7.1.;%
Zhang, Lipton, Li and Smola, \textit{Dive into Deep Learning}, \S3.4.;%
Shalev-Shwartz and Ben-David, \textit{Understanding Machine Learning}, Ch. 15.;%
Murphy, \textit{Probabilistic Machine Learning: An Introduction}, Ch. 5.%
}

\newcommand{\AMLUnitIIISources}{%
Zhang, Lipton, Li and Smola, \textit{Dive into Deep Learning}, Ch. 12 (Optimization Algorithms).;%
Deisenroth, Faisal and Ong, \textit{Mathematics for Machine Learning}, Ch. 7.;%
Kingma and Ba, ``Adam: A Method for Stochastic Optimization,'' ICLR 2015.;%
Ruder, ``An overview of gradient descent optimization algorithms,'' arXiv:1609.04747, 2016.%
}

\newcommand{\AMLUnitIVSources}{%
Han, Kamber and Pei, \textit{Data Mining: Concepts and Techniques} (3e), Ch. 3.;%
M\"uller and Guido, \textit{Introduction to Machine Learning with Python}, Ch. 3--4.;%
James, Witten, Hastie, Tibshirani and Taylor, \textit{An Introduction to Statistical Learning with Python}, Ch. 5--6, 12.;%
Scikit-learn documentation (preprocessing, model selection), accessed 2026-08-04.%
}

\newcommand{\AMLUnitVSources}{%
James et al., \textit{An Introduction to Statistical Learning with Python}, Ch. 3--6.;%
Bishop, \textit{Pattern Recognition and Machine Learning}, Ch. 3.;%
Deisenroth, Faisal and Ong, \textit{Mathematics for Machine Learning}, Ch. 9.;%
M\"uller and Guido, \textit{Introduction to Machine Learning with Python}, \S2.3.%
}

\newcommand{\AMLUnitVISources}{%
Han, Kamber and Pei, \textit{Data Mining: Concepts and Techniques} (3e), Ch. 8--9.;%
James et al., \textit{An Introduction to Statistical Learning with Python}, Ch. 8--9.;%
Bishop, \textit{Pattern Recognition and Machine Learning}, Ch. 4--5, 7, 14.;%
Zhang et al., \textit{Dive into Deep Learning}, Ch. 5.%
}

\newcommand{\AMLUnitVIISources}{%
Han, Kamber and Pei, \textit{Data Mining: Concepts and Techniques} (3e), Ch. 10.;%
James et al., \textit{An Introduction to Statistical Learning with Python}, Ch. 12.;%
M\"uller and Guido, \textit{Introduction to Machine Learning with Python}, \S3.5.;%
Ester, Kriegel, Sander and Xu, ``A density-based algorithm for discovering clusters (DBSCAN),'' KDD 1996.%
}

\newcommand{\AMLLabSources}{%
M\"uller and Guido, \textit{Introduction to Machine Learning with Python}, companion notebooks.;%
McKinney, \textit{Python for Data Analysis} (3e), O'Reilly, 2022.;%
Scikit-learn user guide, accessed 2026-08-04.;%
Matplotlib and Seaborn documentation, accessed 2026-08-04.%
}
%
      }%
    }%
  }%
}
%
  }%
}

\usepackage{fancyhdr}
\usepackage{tcolorbox}
\tcbuselibrary{skins,breakable}

\usepackage{array}
\usepackage[htt]{hyphenat}   % allow \texttt{...} to hyphenate, else long
                             % identifiers overflow the measure

\hypersetup{colorlinks=true, linkcolor=blue, urlcolor=blue, breaklinks=true}
\setlist{itemsep=0.35em, topsep=0.35em}
\sloppy
\emergencystretch=2em

\providecommand{\CourseCode}{CSAI2017P: Applied Machine Learning (Lab)}
\providecommand{\AssignmentName}{Lab Assignment}

\pagestyle{fancy}
\fancyhf{}
\lhead{\small\CourseCode}
\rhead{\small\AssignmentName}
\cfoot{\thepage}
\setlength{\headheight}{14pt}

% Per-question mark tag, right aligned.
\newcommand{\QMarks}[1]{\hfill \textbf{[#1 marks]}}

% Title block for a brief.
\newcommand{\LAtitle}[4]{%
  \begin{center}
    {\LARGE \textbf{#1}}\\[0.25em]
    {\large \CourseCode}\\[0.25em]
    \normalsize #2
  \end{center}
  \noindent
  \textbf{Instructor:} Dr.\ Tofik Ali \hfill \textbf{Due:} #3\\
  \textbf{Student Name:} \rule{5cm}{0.4pt} \hfill \textbf{SAP ID:} \rule{3.5cm}{0.4pt}\\
  \textbf{Batch:} \rule{2.5cm}{0.4pt} \hfill \textbf{Lab quiz:} #4
  \vspace{0.6em}\hrule\vspace{0.8em}}

% Title block for a marking scheme (no student fields).
\newcommand{\MStitle}[3]{%
  \begin{center}
    {\LARGE \textbf{#1}}\\[0.25em]
    {\large \CourseCode}\\[0.25em]
    \normalsize #2\\[0.4em]
    {\color{red}\textbf{MARKING SCHEME --- NOT FOR CIRCULATION}}
  \end{center}
  \noindent \textbf{Instructor:} Dr.\ Tofik Ali \hfill \textbf{Assignment due:} #3
  \vspace{0.6em}\hrule\vspace{0.8em}}

\newtcolorbox{infobox}[1][Note]{
  breakable, enhanced, colback=gray!6, colframe=gray!55,
  fonttitle=\bfseries\small, title={#1}, boxrule=0.7pt, arc=2pt,
  left=6pt,right=6pt,top=4pt,bottom=4pt}

\newtcolorbox{answerbox}[1][Expected answer]{
  breakable, enhanced, colback=green!4, colframe=green!35!black,
  fonttitle=\bfseries\small, title={#1}, boxrule=0.7pt, arc=2pt,
  left=6pt,right=6pt,top=4pt,bottom=4pt}


% Likelihood tags.
\newcommand{\LA}{{\footnotesize\textbf{[$\star\star\star$]}}}
\newcommand{\LB}{{\footnotesize\textbf{[$\star\star$]}}}
\newcommand{\LC}{{\footnotesize\textbf{[$\star$]}}}
\newcommand{\NUM}{{\footnotesize\textbf{[NUM]}}}
\newcommand{\Q}[1]{\item[\textbf{#1}]}

\setlist[enumerate]{itemsep=0.45em, topsep=0.4em}

\begin{document}

\begin{center}
  {\LARGE \textbf{Mid-Semester Question Bank}}\\[0.3em]
  {\large Units I--IV and Unit V (part) \textbullet\ Lecture packages L1--L10}\\[0.25em]
  {\large \CourseCode}\\[0.4em]
  \normalsize Dr.\ Tofik Ali \textbullet\ Autumn 2026 \textbullet\ 314 questions, 63 of them numerical
\end{center}
\vspace{0.4em}\hrule\vspace{0.8em}

\section*{How to use this}

This is not homework. Nothing in here is marked, nothing is collected, and you
may work through it in any order. It exists because you deserve to know what
the mid-semester syllabus means in practice rather than in syllabus headings.

Every question below is one that could reasonably be set on your paper. They
are drawn from the seven lecture packages you have been taught --- the worked
examples, the key ideas, the cautions and the practice questions --- with the
numbers changed so that you have to do the work rather than recall the answer.

\begin{infobox}[Three tags, and what they mean]
\LA\ \textbf{Near certain.} This was a board item, a predict-first reveal, or a
named examinable skill. If it does not appear on the paper in some form, I will
be mildly surprised.

\LB\ \textbf{Likely.} A worked example or a ``key idea'' box from the notes.

\LC\ \textbf{Possible.} Taught, examinable, but peripheral.

\NUM\ \textbf{Numerical.} Compute it by hand, on paper, showing every step.
These are the questions that separate scripts, and they are the ones students
skip in revision because reading feels like progress and arithmetic feels like
work. Do them.
\end{infobox}

\begin{warnbox}[Do the numericals with a pen]
There are \textbf{63} \NUM\ questions here. Reading a worked solution produces
the feeling of understanding and almost none of the substance. You will find
out in the examination hall which of the two you have. Work each one on paper,
\emph{then} check it against the key.

The answer key is released \textbf{after AS1}. That is deliberate.
\end{warnbox}

\section*{The paper you are preparing for}

\begin{center}\begin{tabular}{@{}llll@{}}\toprule
\textbf{Section} & \textbf{Questions} & \textbf{Rule} & \textbf{Marks}\\
\midrule
A & 5 short answers & all compulsory & $5 \times 2 = 10$\\
B & 4 medium questions & attempt any 3 & $3 \times 5 = 15$\\
C & 1 long question & compulsory & $1 \times 10 = 10$\\
\midrule
\multicolumn{3}{@{}l}{\textbf{Total raw}} & \textbf{35}\\
\multicolumn{3}{@{}l}{Scaled to the mid-semester weight} & \textbf{20}\\
\multicolumn{3}{@{}l}{Duration} & \textbf{90 min}\\
\bottomrule\end{tabular}\end{center}

\begin{warnbox}[Neither the syllabus nor the pattern is confirmed yet --- read this]
The examination cell has not yet issued the mid-semester notice. Both the
paper pattern above \emph{and} the exact unit coverage are \textbf{expectations,
not announcements}, and I will tell you the moment either is confirmed.

\textbf{Why Units I--IV plus Unit V to L10 is the working assumption.} A
mid-semester paper can only examine what has been taught by the time you sit
it. The calendar delivers lecture packages \textbf{L1--L10} before the
mid-semester point --- Units I to IV in full, and Unit V as far as logistic
regression. That is where the paper is expected to come from, and it is what
this bank covers. \textbf{Unit V stops at L10:} multiple linear regression
(L11) and regularization (L12) are not here and are not expected.

\textbf{What to do if it turns out otherwise.} This bank is organised strictly
by unit, so it survives either kind of surprise:

\begin{itemize}[leftmargin=*, itemsep=0.15em]
  \item \textbf{Narrower} (say it stops at Unit IV, or at Unit III): work the
    parts that are in and skip the rest. Nothing else changes; the questions you
    skip become end-semester revision, which you will need anyway.
  \item \textbf{Wider} (reaching past L10 into L11 or L12): everything here
    still applies --- nothing in Units I--V(pt) stops being examinable --- and I
    will issue a supplement rather than reissuing this document.
  \item \textbf{A different section pattern}: only the two mock papers are
    affected, and only in their shape. The questions inside them stay valid,
    because a 5-mark question does not become wrong when it is worth 6.
\end{itemize}

None of these is a reason to delay starting. The overlap between any plausible
mid-semester syllabus and this bank is large, and every question in it is
end-semester material too.
\end{warnbox}

\section*{Coverage map}

\begin{center}\footnotesize\begin{tabular}{@{}p{0.9cm}p{5.4cm}p{1.0cm}p{2.4cm}p{3.0cm}@{}}\toprule
\textbf{Unit} & \textbf{What it covers} & \textbf{Lect.} & \textbf{Questions here} & \textbf{Weight to expect}\\
\midrule
I & ML types, the Python stack, generalisation, overfitting, leakage & L1 & 24 (U1.1--U1.24) & light --- at most 1 short, or a part of one\\
II & MSE, MAE, RMSE, Huber, cross-entropy, hinge, triplet & L2 & 38 (U2.1--U2.38) & light--moderate --- expect 1 short, possibly 1 medium\\
III & GD, SGD, mini-batch, momentum, Adagrad, RMSprop, Adadelta, Adam & L3, L4 & 52 (U3.1--U3.52) & moderate --- expect 1 short and up to 1 medium\\
IV & Cleaning, outliers, transformation, features, encoding, PCA, selection, CV, imbalance & L5, L6, L7 & 96 (U4.1--U4.96) & \textbf{heavy} --- expect 1--2 short and 1 medium, and it is a likely source for the long question\\
V\,(pt) & Least squares, MLE, $R^2$, adequacy, over-fitting, logistic regression & L8, L9, L10 & 84 (U5.1--U5.84) & \textbf{heavy} --- expect 1 short and 1--2 medium, and it is the other likely source for the long question\\
--- & Two full mock papers (B and C) & --- & 20 & ---\\
\bottomrule\end{tabular}\end{center}

\noindent\small The ``weight to expect'' column is \emph{indicative}, not a
prediction of the actual paper --- the entries deliberately overlap, because
which unit supplies which question is not knowable in advance. What it does say
is reliable: weight follows \textbf{teaching time}, not unit numbers. Units IV
and V are three lecture packages each and will dominate; Units I and II are one
package each and will not carry much. Budget your revision that way. If the
coverage is announced differently, this column is the only part of the map that
changes.\normalsize

\vspace{0.6em}\hrule

% =====================================================================
\section*{Unit I --- Introduction and the Python ML stack \quad (L1)}

\subsection*{§A --- Two-mark questions}

\begin{enumerate}[leftmargin=2.2cm]
\Q{U1.1}\LA\ State Mitchell's definition of machine learning, naming $T$, $P$
  and $E$. Then fill in all three for a system that predicts whether a student
  will pass a course.
\Q{U1.2}\LA\ Classify each as supervised, unsupervised or reinforcement
  learning, and name the training signal: (a) forecasting monthly rainfall from
  50 years of records; (b) grouping 20{,}000 news articles by topic with no
  labels; (c) training an agent to play Pong from the score alone.
\Q{U1.3}\LB\ State the containment relationship between artificial
  intelligence, machine learning and deep learning. Give one task where deep
  learning is \emph{not} the right choice, with a reason.
\Q{U1.4}\LB\ Give three situations in which you should \emph{not} use machine
  learning, one line each.
\Q{U1.5}\LA\ Name the job of NumPy, Pandas and scikit-learn in one line each.
\Q{U1.6}\LA\ Explain what \texttt{fit}, \texttt{transform} and \texttt{predict}
  each do. State why \texttt{fit\_transform} must never be called on test data.
\Q{U1.7}\LB\ What are the ``first four calls, every time'' on a new DataFrame,
  and what does each tell you?
\Q{U1.8}\LC\ State the NumPy axis mnemonic. For \texttt{a} of shape $(3,4)$,
  give the shape of \texttt{a.sum(axis=0)} and \texttt{a.sum(axis=1)}.
\Q{U1.9}\LB\ Define overfitting in terms of two numbers. State the single
  diagnostic that reveals it.
\Q{U1.10}\LA\ Define data leakage in one sentence. Give two examples.
\Q{U1.11}\LB\ Why is a three-way train/validation/test split needed once you
  start comparing models? Answer in three lines.
\Q{U1.12}\LC\ A classmate reports ``accuracy 95.33\%'' from a single 25\% test
  split of 150 rows. State what is wrong with the way that number is
  \emph{written}, and what should replace it.
\Q{U1.13}\LB\ Name three model families that need feature scaling and two that
  do not, with the reason each way.
\end{enumerate}

\subsection*{§B --- Five-mark questions}

\begin{enumerate}[leftmargin=2.2cm]
\Q{U1.14}\LA\NUM\ Five training flowers are recorded as (petal length, petal
  width) in cm with their species:
  \begin{center}\begin{tabular}{@{}lll@{}}\toprule
  $(1.5,\,0.3)$ setosa & $(1.4,\,0.2)$ setosa & $(4.8,\,1.5)$ versicolor\\
  $(4.2,\,1.2)$ versicolor & $(5.9,\,2.4)$ virginica & \\
  \bottomrule\end{tabular}\end{center}
  A new flower measures $(4.6,\,1.4)$.
  \begin{enumerate}[label=(\alph*)]
    \item Compute the Euclidean distance to each of the five, to four decimals.
    \item Give the 1-NN prediction.
    \item Give the 3-NN prediction and the vote.
    \item State what would change if petal width had been recorded in
      millimetres instead of centimetres, and why.
  \end{enumerate}
\Q{U1.15}\LA\NUM\ A model predicts $\hat y = (4.5,\,2.5,\,7.0,\,5.0)$ against
  true values $y = (5,\,2,\,8,\,4)$.
  \begin{enumerate}[label=(\alph*)]
    \item Compute MSE, MAE and RMSE.
    \item One prediction is now changed: $\hat y_3$ becomes $1.0$. Recompute
      all three.
    \item State the factor by which each metric grew, and explain the ordering.
  \end{enumerate}
\Q{U1.16}\LB\NUM\ A dataset has 569 rows. State (a) the size of a 20\% test
  split; (b) the training and validation fold sizes under 10-fold
  cross-validation; (c) how many model fits 10-fold CV costs; (d) the number of
  row-fits, using $n(k-1)$.
\Q{U1.17}\LA\ The table below was measured on 30 points from a gentle curve
  with 40\% held out.
  \begin{center}\begin{tabular}{@{}lrr@{}}\toprule
  \textbf{Polynomial degree} & \textbf{Train RMSE} & \textbf{Test RMSE}\\
  \midrule
  1 & 1.32 & 1.60\\ 2 & 0.89 & 1.14\\ 3 & 0.86 & 1.10\\ 12 & 0.57 & 40.52\\
  \bottomrule\end{tabular}\end{center}
  \begin{enumerate}[label=(\alph*)]
    \item Which degree would you deploy, and why?
    \item Describe what is happening at degree 12 in terms of the two columns.
    \item Degree 12 has the lowest training error of the four. Explain in three
      lines why that is not an argument in its favour.
    \item Sketch the shape of both curves against degree and label the region
      called ``overfitting''.
  \end{enumerate}
\Q{U1.18}\LA\ A student runs \texttt{SelectKBest(f\_classif, k=20)} on all
  100 rows of a dataset of 5000 pure-noise features with coin-flip labels, then
  cross-validates a classifier on the 20 selected columns and reports 0.81.
  \begin{enumerate}[label=(\alph*)]
    \item State the accuracy an honest procedure would report, and why.
    \item Explain the mechanism producing 0.81 in five lines.
    \item State the one-line code change that fixes it.
  \end{enumerate}
\Q{U1.19}\LB\ Write out, in twelve lines or fewer, a complete scikit-learn
  workflow that loads a dataset, splits it with stratification and a fixed
  seed, scales the features inside a pipeline, fits a $k$-NN classifier and
  reports test accuracy. Annotate the two lines that prevent leakage.
\Q{U1.20}\LB\ Vectorised NumPy was measured at $1.85$\,ms against $162.9$\,ms
  for a Python loop over $10^6$ elements. Explain where the $88\times$ comes
  from, and name one situation where the loop would be the better choice.
\end{enumerate}

\subsection*{§C --- Ten-mark questions}

\begin{enumerate}[leftmargin=2.2cm]
\Q{U1.21}\LB\ ``The goal of machine learning is never to fit the data; it is to
  generalise.'' Develop this into a full answer covering: what a training score
  measures; why a held-out set is needed; what cross-validation adds over a
  single split; the three-way split and what each part is for; and two distinct
  mechanisms by which a reported score can be optimistically biased. Use a
  worked numerical illustration for at least one point.
\Q{U1.22}\LB\ You are handed a labelled dataset and asked to ``build a model''.
  Set out the complete workflow from receiving the file to reporting a number
  you would defend, naming the scikit-learn object used at every stage, and
  stating at each stage what could go wrong and how the pipeline prevents it.
\Q{U1.23}\LC\ Compare a rule-based spam filter with a learned one across:
  how each is built; how each adapts to new spam; how each is debugged; how
  each is explained to a regulator; and what happens to each when the
  distribution of email changes. Conclude with the conditions under which you
  would choose each.
\Q{U1.24}\LC\ Explain the estimator/transformer API as a design decision.
  Why does scikit-learn insist that every object implements the same three or
  four methods? Discuss what this buys for pipelines, cross-validation and grid
  search, and give a concrete example of a bug the design makes impossible.
\end{enumerate}

\vspace{0.4em}\hrule

% =====================================================================
\section*{Unit II --- Loss functions \quad (L2)}

\subsection*{§A --- Two-mark questions}

\begin{enumerate}[leftmargin=2.2cm]
\Q{U2.1}\LA\ Write the formulae for MSE, MAE and RMSE. State which is most
  sensitive to outliers and justify it \emph{from the formula}.
\Q{U2.2}\LA\ Distinguish a loss, a cost (objective) and a metric. Give one
  example of a problem where the loss and the metric deliberately differ.
\Q{U2.3}\LA\ State the two properties a function must have to serve as a loss.
\Q{U2.4}\LA\ Which constant minimises $\sum_i (y_i - c)^2$? Which minimises
  $\sum_i |y_i - c|$? Which minimises the 0--1 loss?
\Q{U2.5}\LA\ Write the Huber loss with parameter $\delta$. State why the second
  branch is $\delta(|r| - \tfrac12\delta)$ rather than $\delta|r|$.
\Q{U2.6}\LB\ State the units of $\delta$ in the Huber loss, and one wrong way
  to choose it.
\Q{U2.7}\LA\ Why can accuracy not be used as a training loss? Give two reasons.
\Q{U2.8}\LA\ Write the binary cross-entropy formula. State what it becomes for
  a single example whose predicted probability on the correct class is $p$.
\Q{U2.9}\LA\ Categorical and sparse categorical cross-entropy: state the label
  format each expects, and whether the two give different numbers.
\Q{U2.10}\LB\ Why does the sparse form exist at all? Give the memory argument
  with a concrete example.
\Q{U2.11}\LB\ Why do libraries clip predicted probabilities to
  $[\varepsilon, 1-\varepsilon]$ before taking logs?
\Q{U2.12}\LA\ Define the margin $m_i$ for a binary classifier with
  $y \in \{-1,+1\}$. Write the hinge loss in terms of it.
\Q{U2.13}\LA\ State the one behavioural difference between hinge loss and
  logistic loss for a point that is already correctly and confidently
  classified. Say what consequence this has for SVMs.
\Q{U2.14}\LB\ Write the triplet loss. Name its three inputs and say what
  $\alpha$ does.
\Q{U2.15}\LB\ Why is random triplet sampling ineffective, and what is the
  remedy called?
\Q{U2.16}\LC\ What is a surrogate loss? Name the two standard surrogates for
  the 0--1 loss.
\Q{U2.17}\LB\ For each, name the loss you would train with: (a) house prices
  with 2\% data-entry errors; (b) 10-class image labels stored as integers;
  (c) a face-embedding model; (d) a two-class medical screen where you need
  calibrated probabilities.
\end{enumerate}

\subsection*{§B --- Five-mark questions}

\begin{enumerate}[leftmargin=2.2cm]
\Q{U2.18}\LA\NUM\ A model gives $\hat y = (3.5,\,5.0,\,4.5,\,8.0,\,5.0)$
  against $y = (3,\,6,\,2,\,9,\,5)$.
  \begin{enumerate}[label=(\alph*)]
    \item Tabulate the residuals, their squares and their absolute values.
    \item Compute MSE, MAE and RMSE.
    \item Compute the fraction of the total squared error, and of the total
      absolute error, contributed by the single worst point.
    \item Comment in three lines on what those two fractions tell you.
  \end{enumerate}
\Q{U2.19}\LA\NUM\ Using the residuals of U2.18, compute the Huber loss at
  $\delta = 1$, $\delta = 1.5$ and $\delta = 3$. For each, state which points
  took the quadratic branch. Explain in three lines what happens to the Huber
  loss as $\delta$ grows large, and what it converges to.
\Q{U2.20}\LA\NUM\ Six exam marks are $20, 22, 24, 25, 27, 28$.
  \begin{enumerate}[label=(\alph*)]
    \item Compute the mean and the median.
    \item The last mark is now recorded as $148$ by mistake. Recompute both.
    \item State which constant minimises the squared error on the corrupted
      set, and which minimises the absolute error.
    \item Explain in four lines why ``MSE fits the mean, MAE fits the median''
      is the same statement as (c).
  \end{enumerate}
\Q{U2.21}\LA\NUM\ Four emails have true labels $y = (1,\,0,\,1,\,0)$ and the
  model predicts $p = (0.8,\,0.3,\,0.6,\,0.1)$ for class 1.
  \begin{enumerate}[label=(\alph*)]
    \item Write down the probability each example assigns to its \emph{correct}
      class.
    \item Compute $-\log$ of each and the binary cross-entropy.
    \item The third email's probability now falls to $0.02$. Recompute the BCE
      and state the increase.
    \item Explain in three lines why one confidently wrong prediction can
      dominate the loss over a large batch.
  \end{enumerate}
\Q{U2.22}\LB\NUM\ A three-class model outputs
  $P = \begin{pmatrix} 0.6 & 0.3 & 0.1\\ 0.2 & 0.7 & 0.1\\ 0.1 & 0.3 & 0.6\end{pmatrix}$
  and the true classes are $0, 1, 2$.
  \begin{enumerate}[label=(\alph*)]
    \item Compute the categorical cross-entropy using one-hot labels.
    \item Compute the sparse categorical cross-entropy using integer labels.
    \item State the relationship between your two answers and explain it.
    \item State the memory saved by the sparse form for $K = 10{,}000$ classes
      and a batch of 256.
  \end{enumerate}
\Q{U2.23}\LA\NUM\ Complete this table and comment on the last column.
  \begin{center}\begin{tabular}{@{}rrr@{}}\toprule
  $p$ on the correct class & $-\log p$ & ratio to the $p=0.99$ row\\
  \midrule
  0.99 & & \\ 0.90 & & \\ 0.50 & & \\ 0.10 & & \\ 0.01 & & \\
  \bottomrule\end{tabular}\end{center}
  State the exact value of the $p = 0.50$ row and why it is a familiar
  constant.
\Q{U2.24}\LA\NUM\ Four points have true labels $y = (+1,\,-1,\,+1,\,-1)$ and
  scores $f = (1.5,\,-0.4,\,0.2,\,-0.9)$.
  \begin{enumerate}[label=(\alph*)]
    \item Compute the margin of each point.
    \item Compute the hinge loss of each and the mean.
    \item Compute the squared hinge loss of each and the mean.
    \item Every point is correctly classified. Three of them still cost
      something. Explain why in four lines, and say what the loss is asking of
      the model that mere correctness does not.
  \end{enumerate}
\Q{U2.25}\LB\NUM\ Compute the triplet loss for each, taking $\alpha$ as given:
  (a) $d(a,p) = 0.9$, $d(a,n) = 1.6$, $\alpha = 1.0$;
  (b) $d(a,p) = 0.4$, $d(a,n) = 2.1$, $\alpha = 1.0$;
  (c) $d(a,p) = 1.2$, $d(a,n) = 1.4$, $\alpha = 0.5$.
  For each, state whether the triplet produces a gradient, and explain what
  case (b) implies about sampling triplets at random.
\Q{U2.26}\LB\ A regression model trained with MSE reports RMSE $1.17$ and a
  model trained with MAE reports MAE $1.00$ on the same data. A classmate
  concludes the second model is better. Give a full rebuttal.
\Q{U2.27}\LB\ Explain why the MAE is harder to optimise with gradient methods
  than the MSE. Your answer must mention the gradient of each, the point
  $r = 0$, and what happens on an even-sized sample.
\Q{U2.28}\LB\ A classifier is being trained on 4 classes. Sketch and explain
  what goes wrong in each case: (a) softmax output with binary cross-entropy;
  (b) integer labels with categorical cross-entropy; (c) sigmoid output on 4
  units with categorical cross-entropy.
\Q{U2.29}\LC\ Explain the family portrait of margin losses: the 0--1 loss and
  three smooth surrogates plotted against the margin. State which are upper
  bounds on the 0--1 loss, where the hinge loss reaches exactly zero, and how
  hinge and squared hinge differ for very negative margins.
\end{enumerate}

\subsection*{§C --- Ten-mark questions}

\begin{enumerate}[leftmargin=2.2cm]
\Q{U2.30}\LA\ ``Choosing a loss is not choosing how to score your model; it is
  choosing what your model will become.'' Defend this claim in full. Cover: the
  constant-prediction result for MSE, MAE and 0--1 loss; a worked numerical
  demonstration on a small contaminated sample; the Huber compromise and the
  role of $\delta$; and a case where the loss you train with should differ from
  the metric you report.
\Q{U2.31}\LA\ Give a complete account of cross-entropy: the binary form, the
  categorical form, the sparse categorical form; what each expects of the
  labels and of the output layer; why the negative log appears; the cost of
  being confidently wrong with numbers; and why cross-entropy rather than
  accuracy is used as the training objective.
\Q{U2.32}\LB\ Compare MSE, MAE and Huber loss across: definition; sensitivity
  to outliers; the constant they fit; differentiability; gradient behaviour;
  hyperparameters; and the situation each is right for. Include one worked
  numerical comparison on a five-point sample of your own construction.
\Q{U2.33}\LC\ Explain margin-based losses in full: the definition of a margin;
  the hinge loss and its geometry; why it produces support vectors; the squared
  hinge; the triplet loss and metric learning; and why triplet mining is
  necessary.
\end{enumerate}

\subsection*{§D --- Statements to judge (2 marks each: true or false, \emph{with} a reason)}

\begin{enumerate}[leftmargin=2.2cm]
\Q{U2.34}\LB\ ``MAE $= 1.0$ is a better result than MSE $= 1.375$ on the same
  predictions.''
\Q{U2.35}\LB\ ``The Huber loss is differentiable everywhere, including at
  $|r| = \delta$.''
\Q{U2.36}\LB\ ``Categorical and sparse categorical cross-entropy are different
  loss functions.''
\Q{U2.37}\LB\ ``A point classified correctly always contributes zero hinge
  loss.''
\Q{U2.38}\LC\ ``RMSE is just MSE in different units, so the two always rank
  models identically.''
\end{enumerate}

\vspace{0.4em}\hrule

% =====================================================================
\section*{Unit III --- Optimiser functions \quad (L3, L4)}

\subsection*{§A --- Two-mark questions}

\begin{enumerate}[leftmargin=2.2cm]
\Q{U3.1}\LA\ Write the gradient descent update rule in vector form. Define
  every symbol.
\Q{U3.2}\LA\ Least squares has a closed-form solution. Give two reasons to use
  gradient descent instead.
\Q{U3.3}\LA\ Define \emph{epoch} and \emph{update}. How many updates per epoch
  does batch gradient descent perform? SGD? Mini-batch with batch size $B$ on
  $n$ rows?
\Q{U3.4}\LA\ State the stability condition on $\eta$ for a quadratic with
  largest curvature $L$. What value of $\eta$ converges in exactly one step,
  and what is that method called?
\Q{U3.5}\LA\ Write the SGD update. State the property of the per-example
  gradient that makes SGD work at all.
\Q{U3.6}\LB\ State the Robbins--Monro conditions on a step-size schedule and
  give a schedule satisfying both.
\Q{U3.7}\LA\ How does the standard deviation of a mini-batch gradient depend on
  $B$? State the practical consequence in one line.
\Q{U3.8}\LA\ Write the momentum update rule. State what $\beta$ controls and
  what happens as $\beta \to 1$.
\Q{U3.9}\LA\ Give the effective step size of momentum in terms of $\eta$ and
  $\beta$. What does this imply when you raise $\beta$ from $0.9$ to $0.99$?
\Q{U3.10}\LB\ Define the condition number $\kappa$. State the convergence rate
  of gradient descent and of momentum in terms of it.
\Q{U3.11}\LB\ Name a situation in which momentum makes optimisation
  \emph{slower}, and say why.
\Q{U3.12}\LA\ Write the Adagrad update. State the one thing $s_t$ accumulates.
\Q{U3.13}\LA\ State Adagrad's flaw on long training runs. Give the rate at
  which its effective learning rate decays.
\Q{U3.14}\LA\ Write the RMSprop update. State the single symbol that differs
  from Adagrad and what it changes.
\Q{U3.15}\LB\ What does RMSprop give up relative to Adagrad?
\Q{U3.16}\LB\ State Adadelta's motivating argument in one sentence. What does
  it remove from the update?
\Q{U3.17}\LB\ ``$\varepsilon$ is Adadelta's hidden learning rate.'' Explain in
  three lines.
\Q{U3.18}\LA\ Write all four Adam equations.
\Q{U3.19}\LA\ Why does Adam need bias correction? Answer in terms of the
  initialisation of $m$ and $v$.
\Q{U3.20}\LA\ Show that Adam's first corrected step is exactly $\pm\eta$,
  whatever the size of $g_1$.
\Q{U3.21}\LB\ Adam is not simply ``better'' than SGD. State what adaptive
  optimisers actually buy you.
\Q{U3.22}\LB\ Name the two distinct reasons a single scalar $\eta$ can fail.
  State which one momentum addresses and which it does not.
\Q{U3.23}\LB\ Give the four signatures of a loss curve and the action each
  calls for.
\Q{U3.24}\LC\ \texttt{SGDRegressor} has no \texttt{batch\_size} argument.
  What does that tell you about its implementation?
\Q{U3.25}\LC\ Name the two conventions for the momentum update and state the
  factor by which they differ.
\end{enumerate}

\subsection*{§B --- Five-mark questions}

\begin{enumerate}[leftmargin=2.2cm]
\Q{U3.26}\LA\NUM\ Minimise $f(w) = (w-4)^2$ by gradient descent from
  $w_0 = 0$ with $\eta = 0.15$.
  \begin{enumerate}[label=(\alph*)]
    \item Write $f'(w)$ and show the update simplifies to $w \leftarrow 0.7w + 1.2$.
    \item Tabulate $k$, $w_k$, $f(w_k)$ and $f'(w_k)$ for $k = 0,1,2,3$ and
      give $w_4$.
    \item State the ratio by which successive gradients fall, and show it
      equals $|1 - 2\eta|$.
    \item Give the closed form $w_k = 4 - 4(1-2\eta)^k$ and verify $w_4$.
  \end{enumerate}
\Q{U3.27}\LA\NUM\ For $f(w) = 3w^2$:
  \begin{enumerate}[label=(\alph*)]
    \item State $f''$ and the range of $\eta$ for which gradient descent
      converges.
    \item Give the $\eta$ that converges in one step.
    \item For $\eta = 0.4$ and $w_0 = 1$, compute $w_1$ to $w_4$ and describe
      the behaviour.
    \item State what a loss curve would look like in case (c), and the first
      thing you would do about it.
  \end{enumerate}
\Q{U3.28}\LA\NUM\ Fit $\hat y = wx$ by minimising
  $J(w) = \frac1n\sum(y_i - wx_i)^2$ on $(1,3)$, $(2,5)$, $(4,9)$, from
  $w_0 = 0$ with $\eta = 0.05$.
  \begin{enumerate}[label=(\alph*)]
    \item Compute the full-batch gradient at $w_0$ and the resulting $w$ after
      one batch update.
    \item Perform one epoch of SGD, visiting the points in the order given,
      using the per-example gradient $2x_i(wx_i - y_i)$. Report $w$ after each
      of the three updates.
    \item Compute the exact optimum $w^\star = \sum x_iy_i / \sum x_i^2$.
    \item Compare (a) and (b) against (c) and comment in three lines.
  \end{enumerate}
\Q{U3.29}\LA\NUM\ Repeat U3.26(b) with momentum, $v_t = \beta v_{t-1} + g_t$
  and $w_t = w_{t-1} - \eta v_t$, using $\beta = 0.9$, $v_0 = 0$, for three
  steps. Tabulate $g_t$, $v_t$ and $w_t$. Compare $w_3$ with plain gradient
  descent's $w_3$ and with the optimum, and explain what you observe.
\Q{U3.30}\LB\NUM\ You tuned $\eta = 0.01$ with $\beta = 0.9$ and now wish to
  use $\beta = 0.99$.
  \begin{enumerate}[label=(\alph*)]
    \item State the effective step size in each case.
    \item Give the rescaling rule and the new $\eta$.
    \item Explain in three lines what happens if you raise $\beta$ without
      rescaling.
  \end{enumerate}
\Q{U3.31}\LA\NUM\ Perform three steps of \textbf{Adagrad} on $f(w) = (w-4)^2$
  from $w_0 = 0$ with $\eta = 0.5$, $s_0 = 0$, $\varepsilon = 10^{-8}$.
  Tabulate $t$, $g_t$, $s_t$, $\sqrt{s_t}$, the step and $w_t$. Then state, with
  numbers, the percentage by which the step shrank against the percentage by
  which the gradient shrank, and explain the difference.
\Q{U3.32}\LA\NUM\ Repeat U3.31 for \textbf{RMSprop} with $\gamma = 0.9$.
  \begin{enumerate}[label=(\alph*)]
    \item Tabulate the same six columns.
    \item Show that the first step equals $\eta/\sqrt{1-\gamma}$ and evaluate
      it.
    \item At which step does $s_t$ first \emph{fall}, and why can Adagrad's
      $s_t$ never do this?
  \end{enumerate}
\Q{U3.33}\LA\NUM\ Perform three steps of \textbf{Adam} on $f(w) = (w-4)^2$
  from $w_0 = 0$ with $\eta = 0.2$, $\beta_1 = 0.9$, $\beta_2 = 0.999$,
  $\varepsilon = 10^{-8}$. Tabulate $t$, $g_t$, $m_t$, $v_t$, $\hat m_t$,
  $\hat v_t$, $\sqrt{\hat v_t}$, the step and $w_t$. Comment on the three step
  sizes.
\Q{U3.34}\LA\NUM\ Repeat U3.33 \emph{without} bias correction.
  \begin{enumerate}[label=(\alph*)]
    \item Report the three uncorrected steps.
    \item Compute the ratio $(1-\beta_1)/\sqrt{1-\beta_2}$ and show it matches
      the first-step discrepancy.
    \item Explain in four lines why the two moments are biased by different
      factors, and why that mismatch — rather than the shrinkage itself — is
      the problem.
  \end{enumerate}
\Q{U3.35}\LB\NUM\ Bias correction matters until $1/(1-\beta^t)$ is within 1\%
  of 1.
  \begin{enumerate}[label=(\alph*)]
    \item Derive the condition on $t$.
    \item Evaluate it for $\beta = 0.9$, $0.99$ and $0.999$.
    \item Comment on what your answer for $\beta_2 = 0.999$ implies for a
      training run of 2000 steps.
  \end{enumerate}
\Q{U3.36}\LB\NUM\ Adadelta's first step is
  $\Delta w_1 = -\dfrac{\sqrt{d_0 + \varepsilon}}{\sqrt{s_1 + \varepsilon}}g_1$
  with $d_0 = 0$ and $s_1 = (1-\rho)g_1^2$. Take $g_1 = -8$, $\rho = 0.9$.
  Compute $\Delta w_1$ for $\varepsilon = 10^{-6}$, $10^{-8}$ and $10^{-12}$,
  and state the relationship you observe. Explain what this means for the claim
  that Adadelta has no learning rate.
\Q{U3.37}\LB\NUM\ Complete the table and answer the question below it.
  \begin{center}\begin{tabular}{@{}rrr@{}}\toprule
  $B$ & relative gradient noise $1/\sqrt B$ & work per update\\
  \midrule
  1 & & \\ 4 & & \\ 16 & & \\ 64 & & \\ 256 & & \\
  \bottomrule\end{tabular}\end{center}
  Doubling $B$ costs twice the work. By what factor does it reduce the noise?
  State the practical conclusion this leads to about the choice of $B$.
\Q{U3.38}\LB\ Complete the optimiser summary table and add one line of comment
  under each row.
  \begin{center}\begin{tabular}{@{}lllll@{}}\toprule
  \textbf{Method} & $B$ & \textbf{Updates/epoch} & \textbf{Noise} & \textbf{Cost/update}\\
  \midrule
  Batch & & & & \\ Mini-batch & & & & \\ Stochastic & & & & \\
  \bottomrule\end{tabular}\end{center}
\Q{U3.39}\LA\ Write from memory the update rules for SGD, momentum, Adagrad,
  RMSprop and Adam, in that order. For each, state in one line what it added
  over its predecessor.
\Q{U3.40}\LB\ ``Every optimiser in this lecture is the line
  $w_i \leftarrow w_i - \frac{\eta}{\sqrt{s_i}}g_i$; they differ only in what
  $s_i$ accumulates.'' Justify this for Adagrad, RMSprop and Adam, stating
  $s_i$ in each case.
\Q{U3.41}\LB\ Give the dimensional-analysis argument for Adadelta. What are the
  units of a plain SGD step, of an Adagrad step, and of a Newton step? Which is
  dimensionally correct, and why is it not used?
\Q{U3.42}\LB\ A model has 20 dense features and 20 very rare ones. Explain why
  a single learning rate cannot serve both, name the optimiser you would reach
  for, and justify it in terms of what that optimiser accumulates.
\Q{U3.43}\LC\ Explain why standardising features is an \emph{optimisation}
  decision, not merely a preprocessing one. Refer to the condition number and
  to the stability limit $2/L$.
\Q{U3.44}\LB\ Give the six-step recipe for training a model with gradient
  methods, from raw features to a final learning-rate decay.
\end{enumerate}

\subsection*{§C --- Ten-mark questions}

\begin{enumerate}[leftmargin=2.2cm]
\Q{U3.45}\LA\ Give a complete account of the learning rate: the update rule;
  the four regimes of $\eta$ with the multiplier $|1-\eta f''|$ in each; the
  stability threshold $2/L$ and its derivation for a quadratic; the four loss
  curve signatures; and how you would find a good $\eta$ in practice. Include a
  worked numerical example on a one-dimensional quadratic.
\Q{U3.46}\LA\ Trace the development from batch gradient descent to Adam as a
  chain of fixes, at each step naming the problem, the fix, the update rule and
  what the fix cost. Cover batch $\to$ SGD $\to$ mini-batch $\to$ momentum
  $\to$ Adagrad $\to$ RMSprop $\to$ Adam.
\Q{U3.47}\LA\ Explain momentum completely: the ravine problem; the heavy-ball
  rule; the exponentially weighted average interpretation and the unrolled
  form; the effective step $\eta/(1-\beta)$; the modified stability ceiling;
  overshoot; the improvement in convergence rate from $\kappa$ to $\sqrt\kappa$;
  and when momentum does not help.
\Q{U3.48}\LA\ Explain Adam in full: the two moments and what each is for; the
  bias-correction derivation for a constant gradient; the first-step result;
  the size of the error if correction is omitted; the defaults and what each
  controls; and the most common misconception about Adam.
\Q{U3.49}\LB\ Compare Adagrad, RMSprop, Adadelta and Adam across: what is
  accumulated; whether it decays; the size of the first step; the
  hyperparameters; the problem each was invented to fix; and the situation each
  is right for. Support with at least two hand-computed steps of any two of
  them.
\Q{U3.50}\LB\ ``The real benefit of an adaptive optimiser is not a better final
  answer but a wider window of learning rates that work.'' Explain and defend
  this, and describe the experiment you would run to demonstrate it.
\end{enumerate}

\subsection*{§D --- Statements to judge (2 marks each)}

\begin{enumerate}[leftmargin=2.2cm]
\Q{U3.51}\LB\ ``Adam is always a better choice than SGD with momentum.''
\Q{U3.52}\LB\ ``Because Adam is adaptive, it does not need a learning-rate
  schedule.''
\end{enumerate}

\vspace{0.4em}\hrule

% =====================================================================
\section*{Unit IV --- Data preprocessing \quad (L5, L6, L7)}

\subsection*{§A --- Two-mark questions}

\begin{enumerate}[leftmargin=2.2cm]
\Q{U4.1}\LA\ Name Han, Kamber and Pei's four preprocessing tasks.
\Q{U4.2}\LA\ Give the probability that a row survives listwise deletion when
  each of $d$ columns is independently missing with probability $p$. Evaluate
  it for $p = 0.05$, $d = 30$.
\Q{U4.3}\LA\ Name three imputation strategies for a numeric column and one for
  a categorical column, with the situation each suits.
\Q{U4.4}\LA\ State the factor by which mean imputation shrinks the standard
  deviation of a column, in terms of the fraction observed.
\Q{U4.5}\LA\ Define MCAR, MAR and MNAR. Give one example of each.
\Q{U4.6}\LB\ For which of the three mechanisms can a multivariate imputer
  recover the bias, and for which can nothing in the data help?
\Q{U4.7}\LB\ What is a missingness indicator, and when is adding one leakage?
\Q{U4.8}\LA\ Write the $z$-score rule, the Tukey fence rule and the modified
  $z$-score rule for outlier detection.
\Q{U4.9}\LA\ Give the breakdown point of each of the three outlier rules.
\Q{U4.10}\LA\ Define \emph{masking}. Explain in two lines why it defeats the
  $z$-score rule.
\Q{U4.11}\LB\ State the largest $|z|$ attainable in a sample of size $n$.
  For what $n$ does the $|z| > 3$ rule become capable of firing at all?
\Q{U4.12}\LB\ Why is $0.6745$ the constant in the modified $z$-score? Why is
  the fence multiplier $1.5$?
\Q{U4.13}\LB\ Distinguish skew from contamination. State the order in which you
  should transform and detect outliers, with a reason.
\Q{U4.14}\LB\ Name four ways to dispose of an identified outlier, and state
  when each is appropriate.
\Q{U4.15}\LC\ Distinguish equal-width from equal-depth binning. Name the three
  smoothing methods.
\Q{U4.16}\LA\ Write the Box--Cox transformation. State its one restriction on
  the data and name the transformation that removes it.
\Q{U4.17}\LB\ Why does the Box--Cox formula subtract 1 and divide by $\lambda$?
\Q{U4.18}\LB\ How is $\lambda$ chosen in practice? State what it is \emph{not}
  chosen to do.
\Q{U4.19}\LA\ Name the four scalers, write each formula and give each output
  range.
\Q{U4.20}\LA\ Which of the four scalers is completely destroyed by a single
  extreme value, and why? Which is the safest?
\Q{U4.21}\LA\ Why does scaling not change a decision tree's predictions?
  Give the argument from the form of a split.
\Q{U4.22}\LB\ Distinguish scaling from transformation. Which one changes
  skewness?
\Q{U4.23}\LA\ Distinguish nominal from ordinal variables. State the encoding
  each requires.
\Q{U4.24}\LA\ What is the dummy variable trap? State what it breaks and what it
  does \emph{not} break.
\Q{U4.25}\LB\ State the default ordering \texttt{OrdinalEncoder} applies, and
  why that is a dangerous default.
\Q{U4.26}\LB\ What does \texttt{handle\_unknown='ignore'} do, and what happens
  without it?
\Q{U4.27}\LB\ Name Han, Kamber and Pei's three kinds of data reduction.
\Q{U4.28}\LA\ State the curse of dimensionality in terms of nearest and
  farthest distances. Name the class of method it damages.
\Q{U4.29}\LA\ Give the four steps of PCA.
\Q{U4.30}\LA\ What does the $j$th eigenvalue of the covariance matrix equal?
  What does the trace equal?
\Q{U4.31}\LA\ Why must you standardise before PCA? Give the failure mode
  concretely.
\Q{U4.32}\LA\ ``PCA reduces dimensionality by keeping the most important
  features.'' State precisely what is wrong with this sentence.
\Q{U4.33}\LB\ PCA has never seen the labels. Give a concrete case where this
  makes it the wrong tool, and name a supervised alternative.
\Q{U4.34}\LA\ Name the three families of feature selection with one method
  from each, and state the cost and risk of each family.
\Q{U4.35}\LB\ The ANOVA $F$ test and mutual information rank features
  differently. State what each measures and construct a feature that one finds
  and the other misses.
\Q{U4.36}\LB\ ``L1 selects, L2 shrinks.'' Explain.
\Q{U4.37}\LA\ State the single rule that determines what may be fitted on the
  full dataset and what may not.
\Q{U4.38}\LA\ Which leaks more: fitting a scaler on all the data, or performing
  feature selection on all the data? Give the reason.
\Q{U4.39}\LA\ Why must a test set be stratified for classification? Name the
  scikit-learn argument.
\Q{U4.40}\LA\ State the default value of \texttt{shuffle} in \texttt{KFold} and
  the failure it causes on a sorted file.
\Q{U4.41}\LA\ What is the ``unit of independence''? Give three examples where
  it is not the row.
\Q{U4.42}\LA\ Name the splitter for each: repeated measurements per patient;
  time-ordered rows; imbalanced classes; imbalanced classes \emph{and} repeated
  measurements.
\Q{U4.43}\LB\ Give the $k$-fold cost model: number of fits, rows per fit, and
  total row-fits.
\Q{U4.44}\LB\ Why is LOOCV rarely worth it? Give both the cost and the variance
  argument.
\Q{U4.45}\LB\ Why is \texttt{cross\_val\_score(...).std()} not the uncertainty
  of the mean?
\Q{U4.46}\LA\ Distinguish the roles of a validation set and a test set. How
  many times may you look at each?
\Q{U4.47}\LB\ Why is \texttt{GridSearchCV.best\_score\_} not a performance
  estimate? What is?
\Q{U4.48}\LA\ Write the formulae for accuracy, precision, recall, $F_1$ and
  balanced accuracy from a confusion matrix.
\Q{U4.49}\LA\ What is the chance baseline of ROC--AUC? Of PR--AUC?
\Q{U4.50}\LA\ Write the SMOTE interpolation formula and define every symbol.
\Q{U4.51}\LB\ Give two situations in which SMOTE is the wrong tool.
\Q{U4.52}\LA\ Where in a cross-validation loop must resampling happen? State
  what goes wrong if it happens earlier.
\Q{U4.53}\LB\ Name a cheaper alternative to resampling that achieves the same
  effect on recall, and state two further advantages it has.
\end{enumerate}

\subsection*{§B --- Five-mark questions}

\begin{enumerate}[leftmargin=2.2cm]
\Q{U4.54}\LA\NUM\ Eight students sat a test; two scripts were lost. The six
  recorded marks are $48, 55, 60, 67, 72, 78$.
  \begin{enumerate}[label=(\alph*)]
    \item Compute the mean, median and standard deviation of the observed
      marks (use $\text{ddof}=1$ for the sd, and state the population sd too).
    \item Impute the two missing marks with the mean. Recompute the mean and
      the population standard deviation of all eight.
    \item Repeat with median imputation.
    \item Compute $\sqrt{m/n}$ and compare it with the ratio of the imputed sd
      to the observed sd. Explain the agreement.
  \end{enumerate}
\Q{U4.55}\LA\NUM\ Complete the survival table for listwise deletion.
  \begin{center}\begin{tabular}{@{}rrrrr@{}}\toprule
  $p$ & $d=5$ & $d=20$ & $d=30$ & $d=50$\\
  \midrule
  0.01 & & & & \\ 0.05 & & & & \\ 0.10 & & & & \\
  \bottomrule\end{tabular}\end{center}
  \begin{enumerate}[label=(\alph*)]
    \item Fill in all twelve cells.
    \item At $p = 0.05$, find the number of columns at which half the rows are
      lost.
    \item At $p = 0.10$ and $d = 50$, how many raw rows must you collect to
      retain 1000 complete ones?
    \item State in three lines when listwise deletion is nevertheless the right
      choice.
  \end{enumerate}
\Q{U4.56}\LA\NUM\ A sensor column reads
  $20,\ 21,\ 19,\ 22,\ 20,\ 21,\ 23,\ 19,\ 20,\ 22,\ 120,\ 125$.
  \begin{enumerate}[label=(\alph*)]
    \item Compute the mean, the sample standard deviation, the median and the
      MAD.
    \item Compute $Q_1$, $Q_3$, the IQR and the $1.5\times$IQR fences.
    \item Score the value 120 under all three rules: $|z| > 3$, the Tukey
      fence, and the modified $z$ at threshold $3.5$. State which rules fire.
    \item Compute the standard deviation of the ten genuine readings alone, and
      the $z$-score the value 120 would then have. Explain the discrepancy with
      (c) and name the effect.
  \end{enumerate}
\Q{U4.57}\LA\NUM\ Complete the table and answer the question.
  \begin{center}\begin{tabular}{@{}rrl@{}}\toprule
  $n$ & largest attainable $|z|$ & can $|z|>3$ ever fire?\\
  \midrule
  5 & & \\ 8 & & \\ 10 & & \\ 11 & & \\ 20 & & \\
  \bottomrule\end{tabular}\end{center}
  Derive the bound $|z_1| \le (n-1)/\sqrt n$. State what it means for a
  practitioner who writes \texttt{if abs(z) > 3} in a script that processes
  batches of eight readings.
\Q{U4.58}\LB\NUM\ Nine prices are $5, 8, 14, 17, 20, 26, 27, 30, 36$.
  \begin{enumerate}[label=(\alph*)]
    \item Partition them into three equal-depth bins and give each bin's mean,
      median and boundaries.
    \item Write the column smoothed by bin means, by bin medians and by bin
      boundaries.
    \item State the bin widths and bin sizes that equal-width binning would
      have produced instead.
    \item State one thing binning gains and one thing it costs.
  \end{enumerate}
\Q{U4.59}\LB\NUM\ Using the sensor column of U4.56, apply each treatment and
  report the resulting mean and standard deviation: (a) leave it alone;
  (b) winsorise at the 10th and 90th percentiles; (c) cap at the IQR fences;
  (d) drop the points flagged by the modified $z$. Rank the four by how close
  they come to the truth (mean $20.7$, sd $1.34$) and justify your ranking.
\Q{U4.60}\LB\NUM\ The values $2, 6, 18, 54$ are in geometric progression.
  \begin{enumerate}[label=(\alph*)]
    \item Compute the gaps between consecutive values.
    \item Compute the natural logarithm of each and the gaps between those.
    \item State the exact value the log gaps take and why.
    \item Explain in three lines why this makes the log the useful ``middle
      rung'' of the ladder of powers.
  \end{enumerate}
\Q{U4.61}\LB\NUM\ Apply the Box--Cox formula $(x^\lambda - 1)/\lambda$ to
  $x = 1, 4, 16, 64$ for $\lambda = 0.5$, $-0.5$ and $0.001$. Compare the
  $\lambda = 0.001$ column with $\log x$ and explain what you observe.
\Q{U4.62}\LA\NUM\ The values $3, 5, 5, 5, 6, 6, 8, 10$.
  \begin{enumerate}[label=(\alph*)]
    \item Compute the mean, population sd, min, max, median, $Q_1$, $Q_3$ and
      IQR.
    \item Transform the column with \texttt{StandardScaler},
      \texttt{MinMaxScaler}, \texttt{RobustScaler} and \texttt{MaxAbsScaler},
      giving all eight values in each case.
    \item State the output range each scaler guarantees.
  \end{enumerate}
\Q{U4.63}\LA\NUM\ A typing error appends the value $400$ to the column of
  U4.62.
  \begin{enumerate}[label=(\alph*)]
    \item Recompute the mean, sd, median and IQR.
    \item For \texttt{StandardScaler}, \texttt{MinMaxScaler} and
      \texttt{RobustScaler}, compute the span the eight \emph{good} values now
      occupy, and the factor by which each has been compressed.
    \item Name the scaler that survives, and explain why in terms of the
      statistics each one uses.
    \item Explain in three lines why \texttt{MinMaxScaler} is described as
      ``the dangerous one'' despite producing output that looks correct.
  \end{enumerate}
\Q{U4.64}\LA\NUM\ Three customers are recorded as (age, annual income in
  rupees): $A = (30,\ 500\,000)$, $B = (55,\ 505\,000)$,
  $C = (31,\ 900\,000)$.
  \begin{enumerate}[label=(\alph*)]
    \item Compute $d(A,B)^2$ and $d(A,C)^2$ on the raw values, showing each
      coordinate's contribution and the percentage income supplies.
    \item Standardise all three columns (population sd) and recompute both
      squared distances.
    \item State $A$'s nearest neighbour before and after.
    \item Explain in four lines what ``choosing not to scale'' actually chooses.
  \end{enumerate}
\Q{U4.65}\LA\NUM\ A nominal colour variable is ordinal-encoded alphabetically:
  blue $=0$, green $=1$, red $=2$. The true group means of the target are
  blue 30, green 70, red 15.
  \begin{enumerate}[label=(\alph*)]
    \item Compute $\bar x$, $\bar y$, $S_{xy}$ and $S_{xx}$.
    \item Fit the least-squares line, giving slope and intercept.
    \item Give the prediction for each colour, the SSE, the SST and $R^2$.
    \item State the $R^2$ a one-hot encoding would achieve and explain the
      difference in four lines.
  \end{enumerate}
\Q{U4.66}\LB\NUM\ A categorical column has 5 levels.
  \begin{enumerate}[label=(\alph*)]
    \item How many columns does one-hot encoding produce, with and without
      \texttt{drop='first'}?
    \item A design matrix has an intercept column and all 5 dummies. State its
      rank and why.
    \item State what rank deficiency breaks, and what it does \emph{not} break.
    \item State when you should use \texttt{drop='first'} and when you should
      not.
  \end{enumerate}
\Q{U4.67}\LB\NUM\ \texttt{PolynomialFeatures} generates $\binom{d+k}{k} - 1$
  features. Evaluate for $(d,k) = (2,2), (10,2), (30,2), (10,3), (100,2)$.
  Comment on what this implies for ``just generate all the interactions''.
\Q{U4.68}\LA\NUM\ Perform PCA by hand on the five points
  $(1,1), (2,3), (3,2), (4,5), (5,4)$.
  \begin{enumerate}[label=(\alph*)]
    \item Compute the mean and the centred matrix.
    \item Compute the $2\times2$ covariance matrix with divisor $n-1$.
    \item Find both eigenvalues and the first eigenvector. Verify the trace.
    \item Give the proportion of variance explained by each component.
    \item Project onto PC1, give the five scores, and verify their variance
      equals $\lambda_1$.
    \item Reconstruct from one component and compute the total squared
      reconstruction error. Verify it equals $(n-1)\lambda_2$.
  \end{enumerate}
\Q{U4.69}\LB\NUM\ A PCA on six standardised features gives eigenvalues
  $4.8,\ 2.4,\ 1.3,\ 0.8,\ 0.4,\ 0.3$.
  \begin{enumerate}[label=(\alph*)]
    \item Give the proportion of variance explained by each component.
    \item Give the cumulative proportions.
    \item State the smallest $k$ reaching 80\%, 90\% and 95\%.
    \item State the relative reconstruction error at each of those three $k$,
      without further computation, and justify why no computation is needed.
  \end{enumerate}
\Q{U4.70}\LA\NUM\ Twelve rows contain three positives, all at the end of the
  file. Four rows are held out as a test set.
  \begin{enumerate}[label=(\alph*)]
    \item If the last four rows are taken, how many positives are in the test
      set and how many in training? State the consequence.
    \item If the rows are shuffled first, compute the probability that the test
      set contains \emph{no} positive, as a ratio of binomial coefficients.
    \item State exactly what \texttt{stratify=y} produces.
    \item State what \texttt{recall} evaluates to on a fold containing no
      positives, and what scikit-learn does about it.
  \end{enumerate}
\Q{U4.71}\LA\NUM\ A screening model on 10{,}000 patients gives
  $\text{TN} = 9900$, $\text{FP} = 60$, $\text{FN} = 30$, $\text{TP} = 10$.
  \begin{enumerate}[label=(\alph*)]
    \item Compute accuracy, precision, recall, $F_1$ and balanced accuracy.
    \item Compute the accuracy of a model that predicts ``negative'' always.
    \item State which of the two models scores higher on accuracy, and what
      that tells you about the metric.
    \item Name the metric you would report to a clinical committee and justify
      it in three lines.
  \end{enumerate}
\Q{U4.72}\LA\NUM\ Complete the cross-validation cost table for $n = 5000$.
  \begin{center}\begin{tabular}{@{}lrrr@{}}\toprule
  \textbf{Scheme} & \textbf{Fits} & \textbf{Rows per fit} & \textbf{Row-fits}\\
  \midrule
  2-fold & & & \\ 5-fold & & & \\ 10-fold & & & \\ LOOCV & & & \\
  \bottomrule\end{tabular}\end{center}
  If one fit on 4000 rows takes $0.4$ seconds and cost is linear in rows, give
  the wall-clock time of each scheme. State which you would use and why.
\Q{U4.73}\LB\NUM\ A dataset has 5000 rows and an 8\% positive rate, and you run
  5-fold stratified cross-validation.
  \begin{enumerate}[label=(\alph*)]
    \item Give the number of rows and positives in each validation fold and in
      each training portion.
    \item The training portion is randomly oversampled to balance. Give the
      resulting row count and class counts, and the percentage of minority rows
      that are duplicates.
    \item Repeat for random undersampling, and state the percentage of majority
      rows discarded.
    \item State what happens to the validation fold in each case, and why.
  \end{enumerate}
\Q{U4.74}\LB\NUM\ Two minority patients are at $A = (2,3)$ and $B = (6,11)$.
  \begin{enumerate}[label=(\alph*)]
    \item Compute the SMOTE synthetic point for $u = 0,\ 0.25,\ 0.5,\ 0.75,\ 1$.
    \item State the geometric object all synthetic points lie on.
    \item State the assumption SMOTE is making, in one sentence.
    \item Give one dataset for which that assumption is clearly false.
  \end{enumerate}
\Q{U4.75}\LB\NUM\ A model reports accuracy $0.91$ on a test set of 44 rows.
  \begin{enumerate}[label=(\alph*)]
    \item Compute the standard error of that accuracy.
    \item State how much one row is worth in accuracy points.
    \item Compute the test-set size needed to halve the standard error.
    \item Explain in three lines why comparing two models on one 44-row test
      set is unsound.
  \end{enumerate}
\Q{U4.76}\LA\ A student oversamples the minority class, then splits 80/20, and
  reports $F_1 = 0.89$. Doing it the other way round gives $F_1 = 0.11$.
  \begin{enumerate}[label=(\alph*)]
    \item Name the error.
    \item Explain the mechanism in five lines, referring to what appears in
      both the training and test sets.
    \item State what fraction of the minority test rows are likely to be exact
      copies of training rows, and why.
    \item State the one-line fix.
  \end{enumerate}
\end{enumerate}

\subsection*{§C --- Ten-mark questions}

\begin{enumerate}[leftmargin=2.2cm]
\Q{U4.77}\LA\ You receive 5000 patient records: 12 numeric features (two with
  about 4\% missing values and visible outliers), 3 categorical features, and a
  binary target that is 8\% positive. Design and defend a complete pipeline.
  Cover: the order of every preprocessing step and why it sits where it does;
  exactly which objects are fitted on training data only; what happens inside
  one fold of 5-fold stratified cross-validation, with fold sizes and positive
  counts; where resampling goes and what breaks if it goes earlier; and which
  single metric you would report, with a justification in terms of the clinical
  cost of a false negative.
\Q{U4.78}\LA\ Give a complete account of missing data: listwise deletion and
  its exponential cost; the three simple imputers; the variance-shrinkage
  result with its derivation; MCAR, MAR and MNAR with an example of each; which
  mechanisms a multivariate imputer can repair; the missingness indicator and
  when it is itself leakage; and a recommendation with the reason it is not the
  choice of imputer that matters most.
\Q{U4.79}\LA\ Give a complete account of outlier detection and treatment: the
  three rules with their formulae; their breakdown points; masking and
  swamping, with a worked numerical demonstration; the $z$-score ceiling and
  its proof; why skew is not contamination and what to do about it; and the
  four disposal options with the circumstances for each. Conclude with the
  distinction between detection and disposal.
\Q{U4.80}\LA\ Explain feature scaling completely: the three mechanisms by which
  units matter; the four scalers with formulae and output ranges; a worked
  numerical demonstration that a nearest neighbour can change; which model
  families are affected and which are provably not, with the argument; what a
  single typo does to each scaler; and why scaling is fitted on the training
  split.
\Q{U4.81}\LA\ Explain PCA completely: the question it asks; the four-step
  recipe; the meaning of the eigenvalues and of the trace; a worked
  by-hand example on five points including the reconstruction error identity;
  how $k$ is chosen and why there is no universal rule; why standardising comes
  first; why PCA is not feature selection; and why being unsupervised can make
  it the wrong tool.
\Q{U4.82}\LA\ Explain cross-validation completely: what a single held-out score
  estimates and how noisy it is; the $k$-fold recipe and its cost model; how
  $k$ trades bias against variance against time; why LOOCV's across-fold
  standard deviation is not what it appears to be; the three failure modes
  requiring stratified, grouped and time-ordered splitters; and how tuning
  destroys the estimate unless nested.
\Q{U4.83}\LA\ Explain imbalanced data completely: why accuracy fails, with a
  worked confusion matrix; the metrics that do not fail and their chance
  baselines; the four available remedies; SMOTE's arithmetic and its
  assumption; the empirical result that resampling improves recall while
  costing precision; the threshold-moving alternative; and the rule about where
  resampling belongs, with the size of the error when it is broken.
\Q{U4.84}\LA\ ``Anything you learn from the data --- a mean, a category list, a
  feature subset, a hyperparameter, a resampled row --- is part of the model,
  and belongs inside the loop.'' Take this as a thesis and defend it. Give at
  least four distinct leaks, ranked by the size of the error each produces, and
  say for each what the correct construction is.
\Q{U4.85}\LB\ Explain feature selection: the three families with two methods
  each; what a filter measures and why choosing one is choosing an assumption;
  the disagreement between the ANOVA $F$ test and mutual information; recursive
  feature elimination and its cost; how L1 differs from L2; why four selectors
  can disagree almost completely; and why ``select features because it is good
  practice'' is bad advice.
\Q{U4.86}\LB\ Explain feature engineering: what it is; the four moves; a worked
  example where one constructed column changes $R^2$ from 0.66 to 1.00; why a
  representation is good or bad only relative to a model, with the binning
  result; and why automatic feature generation is not a substitute for a
  reason.
\end{enumerate}

\subsection*{§D --- Statements to judge (2 marks each)}

\begin{enumerate}[leftmargin=2.2cm]
\Q{U4.87}\LA\ ``The $z$-score rule detects outliers more than three standard
  deviations from the mean.''
\Q{U4.88}\LA\ ``PCA reduces dimensionality by keeping the most important
  features.''
\Q{U4.89}\LB\ ``Scaling the inputs to a decision tree improves its accuracy.''
\Q{U4.90}\LB\ ``Mean imputation is safe because it does not change the mean.''
\Q{U4.91}\LB\ ``One-hot encoding a variable with 257 levels is fine; the model
  will ignore the useless columns.''
\Q{U4.92}\LA\ ``Fitting the scaler on the whole dataset is a serious leak that
  will badly inflate your score.''
\Q{U4.93}\LB\ ``LOOCV is the most thorough form of cross-validation and should
  be used whenever it is affordable.''
\Q{U4.94}\LA\ ``SMOTE generates new data.''
\Q{U4.95}\LB\ ``A model with 99\% accuracy on a 1\%-positive problem is
  performing well.''
\Q{U4.96}\LC\ ``\texttt{random\_state=42} is a formality.''
\end{enumerate}

\vspace{0.4em}\hrule

% =====================================================================
\section*{Unit V (part) --- Regression \quad (L8, L9, L10)}

\begin{infobox}[This unit is covered only as far as L10]
Unit V runs from L8 to L12. The mid-semester paper reaches only \textbf{L8, L9
and L10} --- regression foundations and least squares, maximum likelihood and
model adequacy, and logistic regression. \textbf{Multiple linear regression
(L11) and regularization (L12) are not in this bank} and are not expected on
the paper.

One consequence worth knowing: L10's separation problem is \emph{fixed} by
regularization, which is L12. So a question here may ask you to \emph{diagnose}
perfect separation and \emph{name} the remedy, but not to work with ridge or
lasso.
\end{infobox}

\subsection*{§A --- Two-mark questions}

\begin{enumerate}[leftmargin=2.2cm]
\Q{U5.1}\LA\ Write the simple linear regression model, naming every symbol.
  State which symbols are Greek and which are Roman, and why the distinction
  matters.
\Q{U5.2}\LA\ ``Linear'' in linear regression means linear in \emph{what}?
  Classify each as linear or not: (a) $y = \beta_0 + \beta_1 x + \beta_2 x^2$;
  (b) $y = \beta_0 + \beta_0\beta_1 x$; (c) $y = \beta_0 + \beta_1 \log x$;
  (d) $y = \beta_0 e^{\beta_1 x}$.
\Q{U5.3}\LA\ State the criterion least squares minimises. Write it out in full,
  and say precisely which deviations are being squared.
\Q{U5.4}\LA\ Write the two normal equations for simple linear regression.
\Q{U5.5}\LA\ Give the least-squares estimators $b_1$ and $b_0$ in terms of
  $S_{xx}$ and $S_{xy}$, and define both.
\Q{U5.6}\LB\ Give three reasons for choosing squared deviations over absolute
  deviations, and one thing you give up.
\Q{U5.7}\LB\ Why vertical deviations rather than perpendicular ones? Name what
  you get if you minimise perpendicular distance instead.
\Q{U5.8}\LA\ State the second-order condition for the least-squares solution.
  Give the Hessian and the exact condition under which the solution exists and
  is unique.
\Q{U5.9}\LB\ Give a dataset on which the least-squares line does not exist, and
  say which quantity is zero.
\Q{U5.10}\LA\ Write the matrix form of the normal equations and its solution.
  State the condition on $X$.
\Q{U5.11}\LA\ State the four properties of the least-squares line. Which one
  requires an intercept?
\Q{U5.12}\LB\ Prove that the residuals sum to zero, from the normal equations.
\Q{U5.13}\LB\ Prove that the line passes through $(\bar x, \bar y)$.
\Q{U5.14}\LA\ Write the ANOVA decomposition of the total sum of squares and
  name each term.
\Q{U5.15}\LA\ Give $\operatorname{Var}(b_1)$. State what a designer of an
  experiment should therefore do when choosing the $x$ values.
\Q{U5.16}\LB\ Define the leverage $h_i$. Give the rule for flagging a
  high-leverage point.
\Q{U5.17}\LB\ Compare the closed form with gradient descent for fitting a
  linear model, across cost, memory, tuning and when each fails.
\Q{U5.18}\LA\ State the probability model behind maximum likelihood estimation
  for regression. What exactly is assumed about $\varepsilon_i$?
\Q{U5.19}\LA\ Write the likelihood and the log-likelihood for the Gaussian
  linear model.
\Q{U5.20}\LA\ State in one sentence the central result linking maximum
  likelihood to least squares.
\Q{U5.21}\LA\ Why does the value of $\sigma^2$ not affect which $\beta_0,
  \beta_1$ maximise the likelihood?
\Q{U5.22}\LB\ Give the maximum likelihood estimate of $\sigma^2$, and the
  unbiased alternative. State the reason for the difference --- \emph{not} by
  restating the formula.
\Q{U5.23}\LB\ Change the noise model from Gaussian to Laplace. What criterion
  does maximum likelihood then minimise?
\Q{U5.24}\LA\ Define $R^2$ two ways. State its range on training data, and why
  that range does not apply to held-out data.
\Q{U5.25}\LA\ Why can $R^2$ never decrease when a predictor is added? Answer in
  one sentence.
\Q{U5.26}\LB\ Write adjusted $R^2$. State what it is and what it is not.
\Q{U5.27}\LA\ ``$R^2 = 0.85$, so the model is good.'' State two distinct things
  wrong with that inference.
\Q{U5.28}\LB\ Name the four assumptions of the linear model. For each, state
  what its violation invalidates.
\Q{U5.29}\LB\ Define over-fitting. State why it is a property of the model
  \emph{and} the data, not the model alone.
\Q{U5.30}\LA\ Why does linear regression fail on a 0/1 target? Give two
  distinct failures.
\Q{U5.31}\LA\ Define odds and log-odds. Give the range of each, and say which
  one $w^\top x + b$ can safely produce.
\Q{U5.32}\LA\ Write the logistic regression model in log-odds form, then invert
  it to get $p(x)$.
\Q{U5.33}\LA\ State three properties of the sigmoid, including its derivative.
\Q{U5.34}\LA\ Write the log-likelihood for logistic regression. State its
  relationship to binary cross-entropy.
\Q{U5.35}\LA\ Write the gradient of the logistic log-likelihood in matrix form.
  State in words what it is measuring.
\Q{U5.36}\LA\ Why is there no closed form for logistic regression? Contrast the
  score equations with the normal equations.
\Q{U5.37}\LB\ Write the Hessian of the logistic log-likelihood and state what
  its definiteness buys you --- and what it does not.
\Q{U5.38}\LA\ Interpret a logistic coefficient $\beta_j$. State exactly what
  $e^{\beta_j}$ is a change in.
\Q{U5.39}\LA\ State the divide-by-four rule and what it bounds.
\Q{U5.40}\LB\ Why is the decision threshold not part of the model? Who should
  choose it?
\Q{U5.41}\LB\ Write the softmax. Show that it reduces to the sigmoid when
  $K = 2$.
\Q{U5.42}\LA\ What is perfect separation? State what happens to the maximum
  likelihood estimate, and name the fix.
\Q{U5.43}\LC\ Why is squared error the wrong loss for logistic regression?
  Give two reasons.
\end{enumerate}

\subsection*{§B --- Five-mark questions}

\begin{enumerate}[leftmargin=2.2cm]

\Q{U5.44}\LA\NUM\ Fit a least-squares line to
  $x = (0, 1, 2, 3, 4)$, $y = (1, 2, 4, 5, 8)$.
  \begin{enumerate}[label=(\alph*)]
    \item Tabulate $x_i - \bar x$, $y_i - \bar y$, $(x_i-\bar x)^2$ and
      $(x_i-\bar x)(y_i-\bar y)$, and give $S_{xx}$, $S_{xy}$, $S_{yy}$.
    \item Compute $b_1$ and $b_0$ and write the fitted line.
    \item Give the fitted values and residuals. Verify $\sum e_i = 0$ and
      $\sum x_i e_i = 0$.
    \item Compute SSE, SSR and SST, and verify the decomposition. Check SSR
      two independent ways.
    \item Predict $y$ at $x = 6$.
  \end{enumerate}
\Q{U5.45}\LA\NUM\ Repeat the fit of U5.44 by the matrix route.
  \begin{enumerate}[label=(\alph*)]
    \item Write the design matrix $X$ and compute $X^\top X$ and $X^\top y$
      from the raw sums.
    \item Compute $\det(X^\top X)$ and verify it equals $n S_{xx}$.
    \item Invert the $2\times2$ matrix and solve for $b$.
    \item State the exact condition under which this route fails, and give a
      five-row dataset on which it does.
  \end{enumerate}
\Q{U5.46}\LA\ Derive the least-squares estimators from a blank page. Your
  answer must contain: the object being minimised; both partial derivatives;
  both normal equations; the substitution that isolates $b_1$; the
  identification of $S_{xx}$ and $S_{xy}$; and the second-order condition.
  \emph{Stopping at the first-order conditions is not a complete answer.}
\Q{U5.47}\LA\NUM\ Four candidate lines are proposed for the data of U5.44.
  \begin{center}\begin{tabular}{@{}llll@{}}\toprule
  A: $\hat y = 0.5 + 1.7x$ & B: $\hat y = 0.6 + 1.7x$ &
  C: $\hat y = 0.6 + 1.8x$ & D: $\hat y = 1.0 + 1.5x$\\
  \bottomrule\end{tabular}\end{center}
  For each, compute SSE, $\sum e_i$ and $\sum x_ie_i$. Identify the
  least-squares line \textbf{using both normal equations}, and explain why one
  of them alone is not enough --- naming the candidate that proves it.
\Q{U5.48}\LB\NUM\ Refit the data of U5.44 \emph{without} an intercept.
  \begin{enumerate}[label=(\alph*)]
    \item Derive the estimator $b_1 = \sum x_iy_i/\sum x_i^2$ from scratch.
    \item Evaluate it.
    \item Give the residuals and their sum.
    \item Property 1 has died. Explain exactly which normal equation is
      missing and why.
    \item Compare the two SSEs and say whether the comparison is fair.
  \end{enumerate}
\Q{U5.49}\LA\NUM\ For the data of U5.44:
  \begin{enumerate}[label=(\alph*)]
    \item Compute the leverage $h_i = 1/n + (x_i-\bar x)^2/S_{xx}$ of every
      point, and verify the leverages sum to 2.
    \item Adding $\delta$ to $y_i$ changes $b_1$ by $\delta(x_i-\bar x)/S_{xx}$.
      Derive this.
    \item Adding 3 to one $y$ value: give the new $b_1$ for each of the five
      choices.
    \item One point cannot move the slope at all. Say which and why.
  \end{enumerate}
\Q{U5.50}\LB\NUM\ $\operatorname{Var}(b_1) = \sigma^2/S_{xx}$ with
  $\sigma^2 = 4$. Three designs, all with the same budget:
  all with $n = 5$: equally spaced $x$ over $[0,8]$; five points clustered on
  $[3,5]$; and three points at $x=0$ with two at $x=8$.
  Compute $S_{xx}$, $\operatorname{Var}(b_1)$ and $\operatorname{sd}(b_1)$ for
  each. Rank them, then state the one thing the best design cannot do.
\Q{U5.51}\LA\ Derive the maximum likelihood estimators for the Gaussian linear
  model from a blank page: the probability model, the one-row density, the
  likelihood, the log-likelihood, and the argument that maximising it is
  identical to minimising SSE for \emph{every} $\sigma^2 > 0$.
\Q{U5.52}\LA\NUM\ Using the fit of U5.44 (SSE $= 1.10$, $n = 5$):
  \begin{enumerate}[label=(\alph*)]
    \item Compute $\hat\sigma^2_{\text{ML}} = \text{SSE}/n$.
    \item Compute the unbiased estimate $s^2 = \text{SSE}/(n-2)$ and the
      standard error of the estimate $s$.
    \item Give the bias factor and the percentage by which the ML estimate
      understates $\sigma^2$ in expectation.
    \item Explain \emph{why} the divisor is $n-2$ --- an answer that restates
      the formula earns nothing.
  \end{enumerate}
\Q{U5.53}\LB\NUM\ For each of $(n, \text{SSE}) = (8, 20)$, $(6, 12)$,
  $(12, 45)$, compute $\hat\sigma^2_{\text{ML}}$, $s^2$, the bias factor and
  the percentage understatement. Then find the smallest $n$ for which the
  understatement falls below 1\%.
\Q{U5.54}\LA\NUM\ For the fit of U5.44:
  \begin{enumerate}[label=(\alph*)]
    \item Compute SS$_{\text{tot}}$, SS$_{\text{res}}$ and SS$_{\text{reg}}$.
    \item Compute $R^2$ from the definition.
    \item Verify $R^2$ two further independent ways.
    \item Compute $r$ and confirm $r^2 = R^2$.
    \item State one thing this $R^2$ does \emph{not} tell you, and what you
      would look at instead.
  \end{enumerate}
\Q{U5.55}\LB\NUM\ With $n = 25$ and $R^2 = 0.55$ held fixed, compute adjusted
  $R^2$ for $p = 1, 5, 10, 15$. Find the smallest $p$ at which adjusted $R^2$
  turns negative, showing the inequality. Explain what a negative adjusted
  $R^2$ is telling you.
\Q{U5.56}\LA\NUM\ On a held-out set, $y = (2, 4, 6)$ and the model predicts
  $(7, 4, 1)$.
  \begin{enumerate}[label=(\alph*)]
    \item Compute SS$_{\text{tot}}$, SS$_{\text{res}}$ and $R^2$.
    \item Compute $R^2$ for a model that predicts the constant 4.
    \item State the lower bound on held-out $R^2$.
    \item Explain in four lines why $R^2 \in [0,1]$ is a training-set
      statement only.
  \end{enumerate}
\Q{U5.57}\LB\NUM\ For $x = (1,2,3,4,5)$, $y = (2,3,6,5,9)$, compute
  $S_{xx}$, $S_{xy}$, $S_{yy}$, the direct slope $S_{xy}/S_{xx}$ and the
  inverse slope $S_{yy}/S_{xy}$. Show that their ratio equals $r^2$ and
  evaluate it. State which line you would report and why, and what goes wrong
  if you invert the direct fit to predict $x$ from $y$.
\Q{U5.58}\LA\ A student reports a degree-3 polynomial on 331 rows with
  training $R^2 = 0.918$, and concludes the model is excellent. The
  cross-validated $R^2$ is $-1212.18$ and the held-out $R^2$ is $-22.04$;
  predicting the training mean scores $-0.0001$.
  \begin{enumerate}[label=(\alph*)]
    \item Name every fault in the reasoning.
    \item State how many parameters a degree-3 expansion of 10 features has,
      and compare it with the row count.
    \item Explain what a held-out $R^2$ of $-22$ means in plain words.
    \item State what should have been reported instead.
  \end{enumerate}
\Q{U5.59}\LA\NUM\ Six tumour sizes $x = (1,2,3,4,5,6)$ mm carry labels
  $y = (0,0,0,1,1,1)$.
  \begin{enumerate}[label=(\alph*)]
    \item Fit a least-squares line. Give $S_{xx}$, $S_{xy}$, the slope and
      the intercept.
    \item Give the six predictions. How many lie outside $[0,1]$?
    \item Give the $x$ at which the prediction crosses $0.5$.
    \item A malignant tumour at $x = 20$ mm is added. Refit and give the new
      boundary. State how far it moved and which original point is now
      misclassified.
    \item Explain the mechanism in three lines, referring to what squared
      error charges for.
  \end{enumerate}
\Q{U5.60}\LA\NUM\ Complete the table and answer what follows.
  \begin{center}\begin{tabular}{@{}rrr@{}}\toprule
  $p$ & odds & log-odds\\
  \midrule
  0.05 & & \\ 0.20 & & \\ 0.50 & & \\ 0.80 & & \\ 0.95 & & \\
  \bottomrule\end{tabular}\end{center}
  State the antisymmetry property you can see in the third column, and say
  what it implies about relabelling the two classes.
\Q{U5.61}\LA\ Derive the logistic regression gradient from a blank page:
  the Bernoulli probability of one row, the likelihood, the log-likelihood,
  the three chain-rule factors, and the cancellation that produces
  $\partial\ell/\partial w_j = \sum_i (y_i - p_i)x_{ij}$.
  \emph{Show the cancellation explicitly.}
\Q{U5.62}\LA\NUM\ Take $X = \begin{pmatrix}1 & -1\\ 1 & 0\\ 1 & 1\\
  1 & 2\end{pmatrix}$, $y = (0,0,1,1)$, $w = (0,0)$, $\eta = 0.4$.
  \begin{enumerate}[label=(\alph*)]
    \item Compute $z$, then $p$, then the log-likelihood. Express the per-row
      negative log-likelihood as a familiar constant.
    \item Compute $p - y$ and the gradient $X^\top(p-y)$.
    \item Perform one gradient step and give the new $w$.
    \item Compute the new $z$, the new $p$ and the new log-likelihood. State
      the improvement.
    \item The intercept component of the gradient is exactly zero. Explain
      why.
  \end{enumerate}
\Q{U5.63}\LA\NUM\ A logistic model has coefficient $\beta = \ln 3$ on a
  feature.
  \begin{enumerate}[label=(\alph*)]
    \item State the odds ratio.
    \item For starting probabilities $0.1$, $0.2$, $0.5$ and $0.8$, compute
      the new probability after a one-unit increase, via the odds.
    \item Give $\Delta p$ in each case and state where it is largest.
    \item Compute the divide-by-four bound and confirm your table respects it.
    \item Correct this sentence: ``a one-unit increase raises the probability
      by 200\%''.
  \end{enumerate}
\Q{U5.64}\LB\NUM\ Compute $\sigma(z)$ and $\sigma(z)(1-\sigma(z))$ for
  $z = -3, -1.5, 0, 1.5, 3$. Verify $\sigma(-z) = 1 - \sigma(z)$ on your
  table, state where the derivative is largest and its value there, and
  explain what that maximum has to do with the divide-by-four rule.
\Q{U5.65}\LA\NUM\ A screening model gives these held-out counts.
  \begin{center}\begin{tabular}{@{}rrrrr@{}}\toprule
  threshold & TN & FP & FN & TP\\
  \midrule
  0.10 & 90 & 20 & 1 & 39\\
  0.25 & 98 & 12 & 3 & 37\\
  0.50 & 105 & 5 & 7 & 33\\
  0.75 & 109 & 1 & 12 & 28\\
  \bottomrule\end{tabular}\end{center}
  \begin{enumerate}[label=(\alph*)]
    \item Compute accuracy, precision, recall and $F_1$ at every threshold.
    \item A missed cancer costs ten times a false alarm. Compute
      $10\,\text{FN} + \text{FP}$ for each row.
    \item State the threshold chosen by accuracy and the threshold chosen by
      cost, and say which you would deploy.
    \item The ROC--AUC is the same for all four rows. Explain why.
  \end{enumerate}
\Q{U5.66}\LB\ A logistic fit returns coefficients $(48.2, -31.7, 0.4)$, an
  accuracy of $1.0000$ and a \texttt{ConvergenceWarning}.
  \begin{enumerate}[label=(\alph*)]
    \item Name the phenomenon.
    \item Explain what happens to the log-likelihood as $\|w\| \to \infty$,
      and why the maximum likelihood estimate does not exist.
    \item State why accuracy cannot detect this.
    \item Name the fix, state what it costs, and say which lecture covers it.
  \end{enumerate}
\Q{U5.67}\LB\ Explain why there is no closed form for logistic regression.
  Write both sets of equations side by side, say precisely what makes one
  solvable and the other not, and name the two iterative methods used instead.
\end{enumerate}

\subsection*{§C --- Ten-mark questions}

\begin{enumerate}[leftmargin=2.2cm]
\Q{U5.68}\LA\ Give a complete account of least squares: the model and its
  assumptions; why squared vertical deviations; the full derivation including
  the second-order condition; the matrix form; the four properties of the
  fitted line; $\operatorname{Var}(b_1) = \sigma^2/S_{xx}$ and its design
  implication; and what least squares will not tell you. Include a worked
  numerical fit.
\Q{U5.69}\LA\ Give a complete account of maximum likelihood for the linear
  model: the probability model; the likelihood and log-likelihood; the central
  equivalence with least squares and why $\sigma^2$ is irrelevant to it; the ML
  estimate of $\sigma^2$, its bias, and the unbiased alternative with the
  degrees-of-freedom reason; and what changes under a Laplace noise model.
\Q{U5.70}\LA\ Give a complete account of $R^2$ and model adequacy: the
  decomposition; the two definitions of $R^2$ and the identity chain in simple
  regression; why $R^2$ cannot decrease with an added predictor; adjusted
  $R^2$ and its limits; negative held-out $R^2$; and why a residual plot must
  accompany every reported $R^2$. Use Anscombe's quartet as the illustration.
\Q{U5.71}\LA\ Give a complete account of over-fitting and its detection: what
  it is; why training error cannot detect it; the degree curve with numbers;
  why it depends on $n$ as well as the model; cross-validation as the
  detector; why one validation split is not enough; and the rule about what
  must be refitted inside every fold.
\Q{U5.72}\LA\ Derive logistic regression completely: why a straight line fails
  on 0/1 labels, with a worked numerical demonstration; odds, log-odds and the
  inversion to the sigmoid; the Bernoulli likelihood and log-likelihood; the
  gradient with the cancellation shown in full; why there is no closed form;
  concavity and what it does and does not guarantee; and one gradient step
  worked by hand.
\Q{U5.73}\LA\ Explain the interpretation and deployment of a logistic model:
  the odds-ratio derivation; why $e^\beta$ is multiplicative in the odds and
  not additive in the probability, with a worked table; the divide-by-four
  rule; why raw coefficients cannot be compared across features; the decision
  threshold as a choice made outside the model; and choosing that threshold
  under asymmetric costs with a worked confusion matrix.
\Q{U5.74}\LB\ Compare least squares and logistic regression across: what is
  being modelled; the loss; whether a closed form exists; the shape of the
  objective; how the parameters are found; how a coefficient is interpreted;
  and what pathology each suffers from. Support with numbers from a worked fit
  of each.
\Q{U5.75}\LB\ ``State the noise distribution, write the likelihood, take
  logarithms, maximise.'' Show that this single recipe produces least squares
  under Gaussian noise, least absolute deviations under Laplace noise, and
  cross-entropy under a Bernoulli model. Do the derivation in each case.
\end{enumerate}

\subsection*{§D --- Statements to judge (2 marks each)}

\begin{enumerate}[leftmargin=2.2cm]
\Q{U5.76}\LA\ ``Least squares is chosen because it is differentiable.''
\Q{U5.77}\LB\ ``The residuals of a least-squares fit always sum to zero.''
\Q{U5.78}\LA\ ``$\hat\sigma^2_{\text{ML}}$ is biased because we divide by $n$
  instead of $n-2$.''
\Q{U5.79}\LA\ ``$R^2$ measures how good the model is.''
\Q{U5.80}\LB\ ``A higher $R^2$ means a better model, so add the predictor.''
\Q{U5.81}\LA\ ``$e^{\beta_j}$ is the increase in the probability.''
\Q{U5.82}\LB\ ``Logistic regression is not a linear model, because the sigmoid
  is not linear.''
\Q{U5.83}\LB\ ``The model reached 100\% accuracy, so the fit is sound.''
\Q{U5.84}\LC\ ``\texttt{LogisticRegression()} in scikit-learn returns the
  maximum likelihood estimate.''
\end{enumerate}

\vspace{0.6em}\hrule\vspace{0.4em}

% =====================================================================
\section*{Mock paper B \quad (90 minutes, 35 marks, scaled to 20)}

\begin{warnbox}[Sit this one properly]
Closed book. No computer. 90 minutes on a timer. Mock paper A is Section D of
Assignment A1; this is the second of three. Do not read the questions before
you start.
\end{warnbox}

\subsection*{Section A --- attempt all \quad (5 $\times$ 2 = 10)}

\begin{enumerate}[label=\textbf{\arabic*.}, leftmargin=*]
\item Write the Huber loss with parameter $\delta$ and state the one property
  that makes it a compromise between MSE and MAE. \QMarks{2}
\item State the stability condition on the learning rate for a quadratic with
  largest curvature $L$, and say what happens at exactly $\eta = 2/L$.
  \QMarks{2}
\item Give the formula for the probability that a row survives listwise
  deletion, and evaluate it for $p = 0.05$ and $d = 20$. \QMarks{2}
\item Name the four scalers and state which one a single extreme value
  destroys, with the reason. \QMarks{2}
\item State the four properties of a least-squares fitted line, and say which
  one requires the model to have an intercept. \QMarks{2}
\end{enumerate}

\subsection*{Section B --- attempt any 3 of 4 \quad (3 $\times$ 5 = 15)}

\begin{enumerate}[label=\textbf{\arabic*.}, leftmargin=*, start=6]
\item Perform three steps of \textbf{RMSprop} on $f(w) = (w-4)^2$ from
  $w_0 = 0$ with $\eta = 0.5$, $\gamma = 0.9$, $\varepsilon = 10^{-8}$.
  Tabulate $t$, $g_t$, $s_t$, $\sqrt{s_t}$, the step and $w_t$. Show that the
  first step equals $\eta/\sqrt{1-\gamma}$, and state at which step $s_t$ first
  falls and why Adagrad's cannot. \QMarks{5}
\item A column reads $20, 21, 19, 22, 20, 21, 23, 19, 20, 22, 120, 125$.
  Compute the mean, sample sd, median, MAD, $Q_1$, $Q_3$, IQR and the
  $1.5\times$IQR fences. Score the value 120 under all three outlier rules and
  state which fire. Name the effect that explains the disagreement. \QMarks{5}
\item Fit a least-squares line to $x = (0,1,2,3,4)$, $y = (1,2,4,5,8)$:
  give $S_{xx}$, $S_{xy}$, $b_1$, $b_0$, the fitted values and the residuals;
  verify $\sum e_i = 0$ and $\sum x_ie_i = 0$; compute SSE, SSR, SST and
  $R^2$, checking SSR a second way. \QMarks{5}
\item Take $X$ with rows $(1,-1), (1,0), (1,1), (1,2)$, $y = (0,0,1,1)$,
  $w = (0,0)$ and $\eta = 0.4$. Compute $p$, the log-likelihood, the gradient
  $X^\top(p-y)$, one gradient step, and the new log-likelihood. Explain why
  the intercept component of the gradient is exactly zero. \QMarks{5}
\end{enumerate}

\subsection*{Section C --- compulsory \quad (10)}

\begin{enumerate}[label=\textbf{\arabic*.}, leftmargin=*, start=10]
\item Trace the development from batch gradient descent to Adam as a chain of
  fixes.
  \begin{enumerate}[label=(\alph*)]
    \item For each of batch GD, SGD, mini-batch, momentum, Adagrad, RMSprop and
      Adam, state the problem it was invented to fix and write its update rule.
      \QMarks{4}
    \item Explain why a single scalar learning rate can fail for two entirely
      distinct reasons, and state which of the two momentum addresses.
      \QMarks{2}
    \item Derive Adam's bias correction for a constant gradient $g$, and show
      that the first corrected step is exactly $\pm\eta$. \QMarks{2}
    \item ``Adam is better than SGD.'' Give the strongest case for this claim
      and the strongest case against it, and state what adaptive optimisers
      actually buy. \QMarks{2}
  \end{enumerate}
\end{enumerate}

\vspace{0.4em}\hrule\vspace{0.4em}

% =====================================================================
\section*{Mock paper C \quad (90 minutes, 35 marks, scaled to 20)}

\subsection*{Section A --- attempt all \quad (5 $\times$ 2 = 10)}

\begin{enumerate}[label=\textbf{\arabic*.}, leftmargin=*]
\item State Mitchell's definition of learning and name $T$, $P$ and $E$ for a
  credit-scoring system. \QMarks{2}
\item Write the binary cross-entropy formula and state what it evaluates to for
  a single example given probability $0.5$ to the correct class. \QMarks{2}
\item Write the momentum update rule and give the effective step size in terms
  of $\eta$ and $\beta$. \QMarks{2}
\item State the factor by which mean imputation shrinks a column's standard
  deviation, and evaluate it when 25\% of values are missing. \QMarks{2}
\item Define odds and log-odds, give the range of each, and state which of the
  two a linear expression $w^\top x + b$ can safely produce. \QMarks{2}
\end{enumerate}

\subsection*{Section B --- attempt any 3 of 4 \quad (3 $\times$ 5 = 15)}

\begin{enumerate}[label=\textbf{\arabic*.}, leftmargin=*, start=6]
\item A model gives $\hat y = (3.5, 5.0, 4.5, 8.0, 5.0)$ against
  $y = (3, 6, 2, 9, 5)$. Compute MSE, MAE, RMSE and the Huber loss at
  $\delta = 1.5$, showing the branch each point takes. Compute the share of the
  total squared error and of the total absolute error carried by the worst
  point, and comment. \QMarks{5}
\item A nominal colour variable is ordinal-encoded blue $=0$, green $=1$,
  red $=2$, with target group means 30, 70 and 15. Compute $S_{xy}$, $S_{xx}$,
  the least-squares slope and intercept, the three predictions, SSE, SST and
  $R^2$. State what one-hot encoding would achieve and explain the difference.
  \QMarks{5}
\item A fit on five rows gives SSE $= 1.10$ with $n = 5$ and
  SS$_{\text{tot}} = 30$. Compute $R^2$, $\hat\sigma^2_{\text{ML}}$, the
  unbiased $s^2$ and $s$. Give the bias factor and the percentage
  understatement, and explain \emph{why} the divisor is $n-2$ --- restating
  the formula earns nothing. \QMarks{5}
\item A logistic model has coefficient $\beta = \ln 3$. State the odds ratio;
  compute the new probability after a one-unit increase from starting
  probabilities $0.1$, $0.2$, $0.5$ and $0.8$, via the odds; state where
  $\Delta p$ is largest; give the divide-by-four bound; and correct the
  sentence ``a one-unit increase raises the probability by 200\%''. \QMarks{5}
\end{enumerate}

\subsection*{Section C --- compulsory \quad (10)}

\begin{enumerate}[label=\textbf{\arabic*.}, leftmargin=*, start=10]
\item Logistic regression, from the failure of the straight line to a deployed
  threshold.
  \begin{enumerate}[label=(\alph*)]
    \item Fit a least-squares line to $x = (1,2,3,4,5,6)$ mm with labels
      $y = (0,0,0,1,1,1)$. Give the slope, the intercept, the six predictions
      and the $0.5$ crossing. State two distinct things wrong with using this
      as a classifier. \QMarks{3}
    \item Derive the sigmoid by inverting the logit link, and state three
      properties of $\sigma$ including its derivative. \QMarks{2}
    \item Derive $\partial\ell/\partial w_j = \sum_i (y_i - p_i)x_{ij}$ from
      the Bernoulli log-likelihood, showing the cancellation explicitly.
      \QMarks{3}
    \item There is no closed form for $w$. Say precisely why, contrasting the
      score equations with the normal equations, and name what concavity does
      and does not guarantee. \QMarks{2}
  \end{enumerate}
\end{enumerate}

\vspace{0.6em}\hrule\vspace{0.5em}

\section*{Eleven ways to lose marks you have earned}

\begin{enumerate}[leftmargin=*, itemsep=0.3em]
\item \textbf{Restating the definition instead of answering.} ``The $z$-score
  rule detects outliers more than three standard deviations from the mean'' is
  the definition. If the question asks why it failed, that sentence earns
  nothing.
\item \textbf{A number with no working.} Full marks need the steps. A correct
  final value alone typically earns 1 of 5.
\item \textbf{Working with no number.} The converse also fails. Finish the
  arithmetic.
\item \textbf{Naming without justifying.} ``Use the IQR rule'' is half an
  answer. \emph{Why} is the other half, and it usually carries the mark.
\item \textbf{Forgetting the units.} An MAE of 1.0 in lakhs of rupees is not an
  MAE of 1.0.
\item \textbf{Answering a different question.} If asked for the effect on
  \emph{precision}, do not write about accuracy.
\item \textbf{One example given twice.} Asked for two examples of leakage,
  ``scaling on all the data'' and ``normalising before the split'' are one.
\item \textbf{Ignoring the mark allocation.} A 2-mark question wants four lines.
  A 10-mark question wants structure, a worked illustration and a conclusion.
  Writing four lines for ten marks is the most common way to lose a grade.
\item \textbf{Skipping the numericals in revision.} They are the most
  predictable marks on the paper and the ones most often left blank.
\item \textbf{Stopping a derivation at the first-order conditions.} Both the
  least-squares derivation and the maximum-likelihood one carry marks for the
  \emph{second-order} check. Setting a derivative to zero finds a stationary
  point; it does not establish that you have found a minimum.
\item \textbf{Saying ``odds'' when you mean ``probability'', or the reverse.}
  $e^{\beta}$ multiplies the \emph{odds}. It is not an increase in probability,
  and no amount of surrounding correct working rescues a sentence that says it
  is.
\end{enumerate}

\vspace{0.5em}\hrule\vspace{0.4em}
\noindent\small\textbf{Answer key} released after AS1. Until then, the lecture
notes for \textbf{L1--L10} carry their own practice questions \emph{with} full
solutions --- work those first if you want something to check yourself
against.\normalsize

\end{document}
